% AUTO-GENERATED by gen_blueprint.py — do not edit by hand
% Regenerate with: python3 gen_blueprint.py

\chapter{L-tromino Tiling of Rectangles and Deficient Rectangles}

\begin{definition}
  \label{lTromino}
  \lean{lTromino}
  \uses{Prototile}
  \leanok
  \proven
  The L-tromino shape, anchored at origin
\end{definition}

\begin{definition}
  \label{lTrominoSet}
  \lean{lTrominoSet}
  \uses{Protoset, lTromino}
  \leanok
  \proven
  The L-tromino shape, anchored at origin -/ def lTromino : Prototile := ⟨[(0, 0), (0, 1), (1, 0)], by decide⟩ /-- A protoset containing just the L-tromino
\end{definition}

\section{Utility Lemmas}

\begin{theorem}
  \label{lTromino_covers_3}
  \lean{lTromino_covers_3}
  \uses{PlacedTile, PlacedTile_cells_card, lTrominoSet}
  \leanok
  \proven
  The L-tromino shape, anchored at origin -/ def lTromino : Prototile := ⟨[(0, 0), (0, 1), (1, 0)], by decide⟩ /-- A protoset containing just the L-tromino -/ def lTrominoSet : Protoset Unit := ⟨fun \_ => lTromino⟩ /- \#\# Utility Lemmas -/ /-- Each L-tromino covers exactly 3 cells
\end{theorem}

\section{Existential Tileability}

\begin{definition}
  \label{LTileable}
  \lean{LTileable}
  \uses{Region, Tileable, lTrominoSet, rectangle}
  \leanok
  \proven
  The L-tromino shape, anchored at origin -/ def lTromino : Prototile := ⟨[(0, 0), (0, 1), (1, 0)], by decide⟩ /-- A protoset containing just the L-tromino -/ def lTrominoSet : Protoset Unit := ⟨fun \_ => lTromino⟩ /- \#\# Utility Lemmas -/ /-- Each L-tromino covers exactly 3 cells -/ theorem lTromino...
\end{definition}

\begin{definition}
  \label{tiling_2x3}
  \lean{tiling_2x3}
  \uses{TileSet, lTrominoSet}
  \leanok
  \proven
  The L-tromino shape, anchored at origin -/ def lTromino : Prototile := ⟨[(0, 0), (0, 1), (1, 0)], by decide⟩ /-- A protoset containing just the L-tromino -/ def lTrominoSet : Protoset Unit := ⟨fun \_ => lTromino⟩ /- \#\# Utility Lemmas -/ /-- Each L-tromino covers exactly 3 cells -/ theorem lTromino...
\end{definition}

\begin{theorem}
  \label{tiling_2x3_valid}
  \lean{tiling_2x3_valid}
  \uses{LTileable, Valid, example, rectangle,
        tiling_2x3}
  \leanok
  \proven
  \((tiling\_2x3\_valid)\)
\end{theorem}

\section{Impossibility Results}

\begin{theorem}
  \label{LTileable_area_div_3}
  \lean{LTileable.area_div_3}
  \uses{LTileable, Region, card_coveredCells, cellsAt,
        example, lTromino_covers_3, rectangle}
  \leanok
  \proven
  The L-tromino shape, anchored at origin -/ def lTromino : Prototile := ⟨[(0, 0), (0, 1), (1, 0)], by decide⟩ /-- A protoset containing just the L-tromino -/ def lTrominoSet : Protoset Unit := ⟨fun \_ => lTromino⟩ /- \#\# Utility Lemmas -/ /-- Each L-tromino covers exactly 3 cells -/ theorem lTromino...
\end{theorem}

\begin{theorem}
  \label{rotateProto_lTromino_x_span}
  \lean{rotateProto_lTromino_x_span}
  \uses{Cell, lTromino, rotateCell, rotateCell90,
        rotateProto}
  \leanok
  \proven
  The L-tromino shape, anchored at origin -/ def lTromino : Prototile := ⟨[(0, 0), (0, 1), (1, 0)], by decide⟩ /-- A protoset containing just the L-tromino -/ def lTrominoSet : Protoset Unit := ⟨fun \_ => lTromino⟩ /- \#\# Utility Lemmas -/ /-- Each L-tromino covers exactly 3 cells -/ theorem lTromino...
\end{theorem}

\begin{theorem}
  \label{lTromino_placed_x_span}
  \lean{lTromino_placed_x_span}
  \uses{Cell, PlacedTile, PlacedTile_cells, lTrominoSet,
        rectangle, rotateProto_lTromino_x_span, translateCell}
  \leanok
  \proven
  The L-tromino shape, anchored at origin -/ def lTromino : Prototile := ⟨[(0, 0), (0, 1), (1, 0)], by decide⟩ /-- A protoset containing just the L-tromino -/ def lTrominoSet : Protoset Unit := ⟨fun \_ => lTromino⟩ /- \#\# Utility Lemmas -/ /-- Each L-tromino covers exactly 3 cells -/ theorem lTromino...
\end{theorem}

\begin{theorem}
  \label{rectangle_1_n_x_eq_0}
  \lean{rectangle_1_n_x_eq_0}
  \uses{Cell, rectangle}
  \leanok
  \proven
  The L-tromino shape, anchored at origin -/ def lTromino : Prototile := ⟨[(0, 0), (0, 1), (1, 0)], by decide⟩ /-- A protoset containing just the L-tromino -/ def lTrominoSet : Protoset Unit := ⟨fun \_ => lTromino⟩ /- \#\# Utility Lemmas -/ /-- Each L-tromino covers exactly 3 cells -/ theorem lTromino...
\end{theorem}

\begin{theorem}
  \label{not_tileable_1_by_n}
  \lean{not_tileable_1_by_n}
  \uses{LTileable, TileSet, coveredCells, lTromino_placed_x_span,
        rectangle, rectangle_1_n_x_eq_0}
  \leanok
  \proven
  The L-tromino shape, anchored at origin -/ def lTromino : Prototile := ⟨[(0, 0), (0, 1), (1, 0)], by decide⟩ /-- A protoset containing just the L-tromino -/ def lTrominoSet : Protoset Unit := ⟨fun \_ => lTromino⟩ /- \#\# Utility Lemmas -/ /-- Each L-tromino covers exactly 3 cells -/ theorem lTromino...
\end{theorem}

\begin{theorem}
  \label{not_tileable_1x3}
  \lean{not_tileable_1x3}
  \uses{LTileable, not_tileable_1_by_n, rectangle}
  \leanok
  \proven
  The L-tromino shape, anchored at origin -/ def lTromino : Prototile := ⟨[(0, 0), (0, 1), (1, 0)], by decide⟩ /-- A protoset containing just the L-tromino -/ def lTrominoSet : Protoset Unit := ⟨fun \_ => lTromino⟩ /- \#\# Utility Lemmas -/ /-- Each L-tromino covers exactly 3 cells -/ theorem lTromino...
\end{theorem}

\begin{definition}
  \label{swapRotation}
  \lean{swapRotation}
  \leanok
  \proven
  The L-tromino shape, anchored at origin -/ def lTromino : Prototile := ⟨[(0, 0), (0, 1), (1, 0)], by decide⟩ /-- A protoset containing just the L-tromino -/ def lTrominoSet : Protoset Unit := ⟨fun \_ => lTromino⟩ /- \#\# Utility Lemmas -/ /-- Each L-tromino covers exactly 3 cells -/ theorem lTromino...
\end{definition}

\begin{theorem}
  \label{swap_rotateProto_lTromino}
  \lean{swap_rotateProto_lTromino}
  \uses{lTromino, rotateProto, swapCell, swapRotation}
  \leanok
  \proven
  The L-tromino shape, anchored at origin -/ def lTromino : Prototile := ⟨[(0, 0), (0, 1), (1, 0)], by decide⟩ /-- A protoset containing just the L-tromino -/ def lTrominoSet : Protoset Unit := ⟨fun \_ => lTromino⟩ /- \#\# Utility Lemmas -/ /-- Each L-tromino covers exactly 3 cells -/ theorem lTromino...
\end{theorem}

\begin{theorem}
  \label{swapCell_translateCell}
  \lean{swapCell_translateCell}
  \uses{Cell, swapCell, translateCell}
  \leanok
  \proven
  The L-tromino shape, anchored at origin -/ def lTromino : Prototile := ⟨[(0, 0), (0, 1), (1, 0)], by decide⟩ /-- A protoset containing just the L-tromino -/ def lTrominoSet : Protoset Unit := ⟨fun \_ => lTromino⟩ /- \#\# Utility Lemmas -/ /-- Each L-tromino covers exactly 3 cells -/ theorem lTromino...
\end{theorem}

\begin{theorem}
  \label{LTileable_swap}
  \lean{LTileable_swap}
  \uses{LTileable, PlacedTile, PlacedTile_cells, Region,
        TileSet, cellsAt, coveredCells, lTromino,
        lTrominoSet, rectangle, rotateProto, swapCell,
        swapCell_injective, swapRegion, swapRotation, swap_rotateProto_lTromino,
        translateCell}
  \leanok
  \proven
  The L-tromino shape, anchored at origin -/ def lTromino : Prototile := ⟨[(0, 0), (0, 1), (1, 0)], by decide⟩ /-- A protoset containing just the L-tromino -/ def lTrominoSet : Protoset Unit := ⟨fun \_ => lTromino⟩ /- \#\# Utility Lemmas -/ /-- Each L-tromino covers exactly 3 cells -/ theorem lTromino...
\end{theorem}

\begin{theorem}
  \label{not_tileable_n_by_1}
  \lean{not_tileable_n_by_1}
  \uses{LTileable, LTileable_swap, not_tileable_1_by_n, rectangle,
        swap_rectangle}
  \leanok
  \proven
  The L-tromino shape, anchored at origin -/ def lTromino : Prototile := ⟨[(0, 0), (0, 1), (1, 0)], by decide⟩ /-- A protoset containing just the L-tromino -/ def lTrominoSet : Protoset Unit := ⟨fun \_ => lTromino⟩ /- \#\# Utility Lemmas -/ /-- Each L-tromino covers exactly 3 cells -/ theorem lTromino...
\end{theorem}

\section{3×odd Impossibility}

\begin{definition}
  \label{rect3x2}
  \lean{rect3x2}
  \uses{Region, rectangle}
  \leanok
  \proven
  The L-tromino shape, anchored at origin -/ def lTromino : Prototile := ⟨[(0, 0), (0, 1), (1, 0)], by decide⟩ /-- A protoset containing just the L-tromino -/ def lTrominoSet : Protoset Unit := ⟨fun \_ => lTromino⟩ /- \#\# Utility Lemmas -/ /-- Each L-tromino covers exactly 3 cells -/ theorem lTromino...
\end{definition}

\section{Key Lemmas for the 3×odd Impossibility Proof}

\begin{theorem}
  \label{lTromino_y_span_2}
  \lean{lTromino_y_span_2}
  \uses{Cell, PlacedTile, PlacedTile_cells, lTromino,
        lTrominoSet, rotateCell, rotateCell90, rotateProto,
        translateCell}
  \leanok
  \proven
  The L-tromino shape, anchored at origin -/ def lTromino : Prototile := ⟨[(0, 0), (0, 1), (1, 0)], by decide⟩ /-- A protoset containing just the L-tromino -/ def lTrominoSet : Protoset Unit := ⟨fun \_ => lTromino⟩ /- \#\# Utility Lemmas -/ /-- Each L-tromino covers exactly 3 cells -/ theorem lTromino...
\end{theorem}

\begin{theorem}
  \label{tile_covering_00_y_bound}
  \lean{tile_covering_00_y_bound}
  \uses{Cell, PlacedTile, lTrominoSet, lTromino_y_span_2}
  \leanok
  \proven
  The L-tromino shape, anchored at origin -/ def lTromino : Prototile := ⟨[(0, 0), (0, 1), (1, 0)], by decide⟩ /-- A protoset containing just the L-tromino -/ def lTrominoSet : Protoset Unit := ⟨fun \_ => lTromino⟩ /- \#\# Utility Lemmas -/ /-- Each L-tromino covers exactly 3 cells -/ theorem lTromino...
\end{theorem}

\begin{theorem}
  \label{tile_covering_20_y_bound}
  \lean{tile_covering_20_y_bound}
  \uses{Cell, PlacedTile, lTrominoSet, lTromino_y_span_2}
  \leanok
  \proven
  The L-tromino shape, anchored at origin -/ def lTromino : Prototile := ⟨[(0, 0), (0, 1), (1, 0)], by decide⟩ /-- A protoset containing just the L-tromino -/ def lTrominoSet : Protoset Unit := ⟨fun \_ => lTromino⟩ /- \#\# Utility Lemmas -/ /-- Each L-tromino covers exactly 3 cells -/ theorem lTromino...
\end{theorem}

\begin{theorem}
  \label{lTromino_x_span_2}
  \lean{lTromino_x_span_2}
  \uses{Cell, PlacedTile, PlacedTile_cells, lTromino,
        lTrominoSet, rect3x2, rectangle, rotateCell,
        rotateCell90, rotateProto, translateCell}
  \leanok
  \proven
  The L-tromino shape, anchored at origin -/ def lTromino : Prototile := ⟨[(0, 0), (0, 1), (1, 0)], by decide⟩ /-- A protoset containing just the L-tromino -/ def lTrominoSet : Protoset Unit := ⟨fun \_ => lTromino⟩ /- \#\# Utility Lemmas -/ /-- Each L-tromino covers exactly 3 cells -/ theorem lTromino...
\end{theorem}

\begin{theorem}
  \label{tile_covering_00_in_3x2}
  \lean{tile_covering_00_in_3x2}
  \uses{PlacedTile, lTrominoSet, mem_rectangle, rect3x2,
        rectangle, tile_covering_00_y_bound}
  \leanok
  \proven
  The L-tromino shape, anchored at origin -/ def lTromino : Prototile := ⟨[(0, 0), (0, 1), (1, 0)], by decide⟩ /-- A protoset containing just the L-tromino -/ def lTrominoSet : Protoset Unit := ⟨fun \_ => lTromino⟩ /- \#\# Utility Lemmas -/ /-- Each L-tromino covers exactly 3 cells -/ theorem lTromino...
\end{theorem}

\begin{theorem}
  \label{tile_covering_20_in_3x2}
  \lean{tile_covering_20_in_3x2}
  \uses{PlacedTile, lTrominoSet, mem_rectangle, rect3x2,
        rectangle, tile_covering_20_y_bound}
  \leanok
  \proven
  The L-tromino shape, anchored at origin -/ def lTromino : Prototile := ⟨[(0, 0), (0, 1), (1, 0)], by decide⟩ /-- A protoset containing just the L-tromino -/ def lTrominoSet : Protoset Unit := ⟨fun \_ => lTromino⟩ /- \#\# Utility Lemmas -/ /-- Each L-tromino covers exactly 3 cells -/ theorem lTromino...
\end{theorem}

\begin{theorem}
  \label{tiles_00_20_disjoint}
  \lean{tiles_00_20_disjoint}
  \uses{TileSet, Valid, lTrominoSet, rect3x2,
        rectangle}
  \leanok
  \proven
  The L-tromino shape, anchored at origin -/ def lTromino : Prototile := ⟨[(0, 0), (0, 1), (1, 0)], by decide⟩ /-- A protoset containing just the L-tromino -/ def lTrominoSet : Protoset Unit := ⟨fun \_ => lTromino⟩ /- \#\# Utility Lemmas -/ /-- Each L-tromino covers exactly 3 cells -/ theorem lTromino...
\end{theorem}

\begin{theorem}
  \label{tiles_00_20_cover_3x2}
  \lean{tiles_00_20_cover_3x2}
  \uses{TileSet, Valid, lTrominoSet, lTromino_covers_3,
        rect3x2, rectangle, tile_covering_00_in_3x2, tile_covering_20_in_3x2,
        tiles_00_20_disjoint}
  \leanok
  \proven
  The L-tromino shape, anchored at origin -/ def lTromino : Prototile := ⟨[(0, 0), (0, 1), (1, 0)], by decide⟩ /-- A protoset containing just the L-tromino -/ def lTrominoSet : Protoset Unit := ⟨fun \_ => lTromino⟩ /- \#\# Utility Lemmas -/ /-- Each L-tromino covers exactly 3 cells -/ theorem lTromino...
\end{theorem}

\begin{theorem}
  \label{remaining_tiles_cover}
  \lean{remaining_tiles_cover}
  \uses{TileSet, Valid, coveredCells, lTrominoSet,
        rect3x2, rectangle, tiles_00_20_cover_3x2}
  \leanok
  \proven
  The L-tromino shape, anchored at origin -/ def lTromino : Prototile := ⟨[(0, 0), (0, 1), (1, 0)], by decide⟩ /-- A protoset containing just the L-tromino -/ def lTrominoSet : Protoset Unit := ⟨fun \_ => lTromino⟩ /- \#\# Utility Lemmas -/ /-- Each L-tromino covers exactly 3 cells -/ theorem lTromino...
\end{theorem}

\begin{theorem}
  \label{rectangle_minus_3x2}
  \lean{rectangle_minus_3x2}
  \uses{mem_rectangle, rect3x2, rectangle, translateCell,
        translateRegion}
  \leanok
  \proven
  The L-tromino shape, anchored at origin -/ def lTromino : Prototile := ⟨[(0, 0), (0, 1), (1, 0)], by decide⟩ /-- A protoset containing just the L-tromino -/ def lTrominoSet : Protoset Unit := ⟨fun \_ => lTromino⟩ /- \#\# Utility Lemmas -/ /-- Each L-tromino covers exactly 3 cells -/ theorem lTromino...
\end{theorem}

\section{Full Rectangle Characterization}

\begin{definition}
  \label{RectTileableConditions}
  \lean{RectTileableConditions}
  \leanok
  \proven
  The L-tromino shape, anchored at origin -/ def lTromino : Prototile := ⟨[(0, 0), (0, 1), (1, 0)], by decide⟩ /-- A protoset containing just the L-tromino -/ def lTrominoSet : Protoset Unit := ⟨fun \_ => lTromino⟩ /- \#\# Utility Lemmas -/ /-- Each L-tromino covers exactly 3 cells -/ theorem lTromino...
\end{definition}

\begin{theorem}
  \label{empty_LTileable}
  \lean{empty_LTileable}
  \uses{LTileable, empty_tileable, lTrominoSet}
  \leanok
  \proven
  The L-tromino shape, anchored at origin -/ def lTromino : Prototile := ⟨[(0, 0), (0, 1), (1, 0)], by decide⟩ /-- A protoset containing just the L-tromino -/ def lTrominoSet : Protoset Unit := ⟨fun \_ => lTromino⟩ /- \#\# Utility Lemmas -/ /-- Each L-tromino covers exactly 3 cells -/ theorem lTromino...
\end{theorem}

\section{Combining Tilings}

\begin{theorem}
  \label{LTileable_translate}
  \lean{LTileable_translate}
  \uses{Cell, LTileable, Region, Tileable_translate,
        lTrominoSet, translateRegion}
  \leanok
  \proven
  The L-tromino shape, anchored at origin -/ def lTromino : Prototile := ⟨[(0, 0), (0, 1), (1, 0)], by decide⟩ /-- A protoset containing just the L-tromino -/ def lTrominoSet : Protoset Unit := ⟨fun \_ => lTromino⟩ /- \#\# Utility Lemmas -/ /-- Each L-tromino covers exactly 3 cells -/ theorem lTromino...
\end{theorem}

\begin{theorem}
  \label{LTileable_translate_iff}
  \lean{LTileable_translate_iff}
  \uses{Cell, LTileable, Region, Tileable_translate_iff,
        lTrominoSet, translate, translateRegion}
  \leanok
  \proven
  The L-tromino shape, anchored at origin -/ def lTromino : Prototile := ⟨[(0, 0), (0, 1), (1, 0)], by decide⟩ /-- A protoset containing just the L-tromino -/ def lTrominoSet : Protoset Unit := ⟨fun \_ => lTromino⟩ /- \#\# Utility Lemmas -/ /-- Each L-tromino covers exactly 3 cells -/ theorem lTromino...
\end{theorem}

\begin{theorem}
  \label{LTileable_of_translate_eq}
  \lean{LTileable_of_translate_eq}
  \uses{Cell, LTileable, Region, Tileable_of_translate_eq,
        lTrominoSet, rectangle, translateRegion}
  \leanok
  \proven
  The L-tromino shape, anchored at origin -/ def lTromino : Prototile := ⟨[(0, 0), (0, 1), (1, 0)], by decide⟩ /-- A protoset containing just the L-tromino -/ def lTrominoSet : Protoset Unit := ⟨fun \_ => lTromino⟩ /- \#\# Utility Lemmas -/ /-- Each L-tromino covers exactly 3 cells -/ theorem lTromino...
\end{theorem}

\section{3×odd Impossibility (requires translation invariance)}

\begin{theorem}
  \label{not_tileable_3_by_odd}
  \lean{not_tileable_3_by_odd}
  \uses{LTileable, LTileable_of_translate_eq, TileSet, cellsAt,
        coveredCells, lTrominoSet, lTromino_x_span_2, mem_rectangle,
        not_tileable_n_by_1, rect3x2, rectangle, rectangle_minus_3x2,
        remaining_tiles_cover}
  \leanok
  \proven
  The L-tromino shape, anchored at origin -/ def lTromino : Prototile := ⟨[(0, 0), (0, 1), (1, 0)], by decide⟩ /-- A protoset containing just the L-tromino -/ def lTrominoSet : Protoset Unit := ⟨fun \_ => lTromino⟩ /- \#\# Utility Lemmas -/ /-- Each L-tromino covers exactly 3 cells -/ theorem lTromino...
\end{theorem}

\begin{theorem}
  \label{not_tileable_odd_by_3}
  \lean{not_tileable_odd_by_3}
  \uses{LTileable, LTileable_swap, not_tileable_3_by_odd, rectangle,
        swap_rectangle}
  \leanok
  \proven
  The L-tromino shape, anchored at origin -/ def lTromino : Prototile := ⟨[(0, 0), (0, 1), (1, 0)], by decide⟩ /-- A protoset containing just the L-tromino -/ def lTrominoSet : Protoset Unit := ⟨fun \_ => lTromino⟩ /- \#\# Utility Lemmas -/ /-- Each L-tromino covers exactly 3 cells -/ theorem lTromino...
\end{theorem}

\begin{theorem}
  \label{rectangle_0_m}
  \lean{rectangle_0_m}
  \uses{rectangle}
  \leanok
  \proven
  The L-tromino shape, anchored at origin -/ def lTromino : Prototile := ⟨[(0, 0), (0, 1), (1, 0)], by decide⟩ /-- A protoset containing just the L-tromino -/ def lTrominoSet : Protoset Unit := ⟨fun \_ => lTromino⟩ /- \#\# Utility Lemmas -/ /-- Each L-tromino covers exactly 3 cells -/ theorem lTromino...
\end{theorem}

\begin{theorem}
  \label{rectangle_n_0}
  \lean{rectangle_n_0}
  \uses{rectangle}
  \leanok
  \proven
  The L-tromino shape, anchored at origin -/ def lTromino : Prototile := ⟨[(0, 0), (0, 1), (1, 0)], by decide⟩ /-- A protoset containing just the L-tromino -/ def lTrominoSet : Protoset Unit := ⟨fun \_ => lTromino⟩ /- \#\# Utility Lemmas -/ /-- Each L-tromino covers exactly 3 cells -/ theorem lTromino...
\end{theorem}

\section{Rectangle Decomposition Lemmas}

\begin{theorem}
  \label{LTileable_swap_rectangle}
  \lean{LTileable_swap_rectangle}
  \uses{LTileable, LTileable_swap, rectangle, swap_rectangle}
  \leanok
  \proven
  The L-tromino shape, anchored at origin -/ def lTromino : Prototile := ⟨[(0, 0), (0, 1), (1, 0)], by decide⟩ /-- A protoset containing just the L-tromino -/ def lTrominoSet : Protoset Unit := ⟨fun \_ => lTromino⟩ /- \#\# Utility Lemmas -/ /-- Each L-tromino covers exactly 3 cells -/ theorem lTromino...
\end{theorem}

\begin{theorem}
  \label{LTileable_horizontal_union}
  \lean{LTileable_horizontal_union}
  \uses{LTileable, Tileable_horizontal_union, lTrominoSet, rectangle}
  \leanok
  \proven
  The L-tromino shape, anchored at origin -/ def lTromino : Prototile := ⟨[(0, 0), (0, 1), (1, 0)], by decide⟩ /-- A protoset containing just the L-tromino -/ def lTrominoSet : Protoset Unit := ⟨fun \_ => lTromino⟩ /- \#\# Utility Lemmas -/ /-- Each L-tromino covers exactly 3 cells -/ theorem lTromino...
\end{theorem}

\begin{theorem}
  \label{LTileable_vertical_union}
  \lean{LTileable_vertical_union}
  \uses{LTileable, Tileable_vertical_union, lTrominoSet, rectangle}
  \leanok
  \proven
  The L-tromino shape, anchored at origin -/ def lTromino : Prototile := ⟨[(0, 0), (0, 1), (1, 0)], by decide⟩ /-- A protoset containing just the L-tromino -/ def lTrominoSet : Protoset Unit := ⟨fun \_ => lTromino⟩ /- \#\# Utility Lemmas -/ /-- Each L-tromino covers exactly 3 cells -/ theorem lTromino...
\end{theorem}

\section{Base Case Tilings}

\begin{theorem}
  \label{tileable_2x3}
  \lean{tileable_2x3}
  \uses{LTileable, rectangle, tiling_2x3, tiling_2x3_valid}
  \leanok
  \proven
  The L-tromino shape, anchored at origin -/ def lTromino : Prototile := ⟨[(0, 0), (0, 1), (1, 0)], by decide⟩ /-- A protoset containing just the L-tromino -/ def lTrominoSet : Protoset Unit := ⟨fun \_ => lTromino⟩ /- \#\# Utility Lemmas -/ /-- Each L-tromino covers exactly 3 cells -/ theorem lTromino...
\end{theorem}

\begin{theorem}
  \label{tileable_3x2}
  \lean{tileable_3x2}
  \uses{LTileable, LTileable_swap_rectangle, rectangle, tileable_2x3}
  \leanok
  \proven
  The L-tromino shape, anchored at origin -/ def lTromino : Prototile := ⟨[(0, 0), (0, 1), (1, 0)], by decide⟩ /-- A protoset containing just the L-tromino -/ def lTrominoSet : Protoset Unit := ⟨fun \_ => lTromino⟩ /- \#\# Utility Lemmas -/ /-- Each L-tromino covers exactly 3 cells -/ theorem lTromino...
\end{theorem}

\begin{theorem}
  \label{tileable_3x6}
  \lean{tileable_3x6}
  \uses{LTileable, LTileable_vertical_union, rectangle, tileable_3x2}
  \leanok
  \proven
  The L-tromino shape, anchored at origin -/ def lTromino : Prototile := ⟨[(0, 0), (0, 1), (1, 0)], by decide⟩ /-- A protoset containing just the L-tromino -/ def lTrominoSet : Protoset Unit := ⟨fun \_ => lTromino⟩ /- \#\# Utility Lemmas -/ /-- Each L-tromino covers exactly 3 cells -/ theorem lTromino...
\end{theorem}

\begin{theorem}
  \label{tileable_2x6}
  \lean{tileable_2x6}
  \uses{LTileable, LTileable_vertical_union, rectangle, tileable_2x3}
  \leanok
  \proven
  The L-tromino shape, anchored at origin -/ def lTromino : Prototile := ⟨[(0, 0), (0, 1), (1, 0)], by decide⟩ /-- A protoset containing just the L-tromino -/ def lTrominoSet : Protoset Unit := ⟨fun \_ => lTromino⟩ /- \#\# Utility Lemmas -/ /-- Each L-tromino covers exactly 3 cells -/ theorem lTromino...
\end{theorem}

\begin{theorem}
  \label{tileable_5x6}
  \lean{tileable_5x6}
  \uses{LTileable, LTileable_horizontal_union, rectangle, tileable_2x6,
        tileable_3x6}
  \leanok
  \proven
  The L-tromino shape, anchored at origin -/ def lTromino : Prototile := ⟨[(0, 0), (0, 1), (1, 0)], by decide⟩ /-- A protoset containing just the L-tromino -/ def lTrominoSet : Protoset Unit := ⟨fun \_ => lTromino⟩ /- \#\# Utility Lemmas -/ /-- Each L-tromino covers exactly 3 cells -/ theorem lTromino...
\end{theorem}

\begin{theorem}
  \label{tileable_6x5}
  \lean{tileable_6x5}
  \uses{LTileable, LTileable_swap_rectangle, rectangle, tileable_5x6}
  \leanok
  \proven
  The L-tromino shape, anchored at origin -/ def lTromino : Prototile := ⟨[(0, 0), (0, 1), (1, 0)], by decide⟩ /-- A protoset containing just the L-tromino -/ def lTrominoSet : Protoset Unit := ⟨fun \_ => lTromino⟩ /- \#\# Utility Lemmas -/ /-- Each L-tromino covers exactly 3 cells -/ theorem lTromino...
\end{theorem}

\begin{theorem}
  \label{tileable_2x9}
  \lean{tileable_2x9}
  \uses{LTileable, LTileable_vertical_union, rectangle, tileable_2x3}
  \leanok
  \proven
  The L-tromino shape, anchored at origin -/ def lTromino : Prototile := ⟨[(0, 0), (0, 1), (1, 0)], by decide⟩ /-- A protoset containing just the L-tromino -/ def lTrominoSet : Protoset Unit := ⟨fun \_ => lTromino⟩ /- \#\# Utility Lemmas -/ /-- Each L-tromino covers exactly 3 cells -/ theorem lTromino...
\end{theorem}

\begin{definition}
  \label{tiling_5x9}
  \lean{tiling_5x9}
  \uses{TileSet, lTrominoSet}
  \leanok
  \proven
  The L-tromino shape, anchored at origin -/ def lTromino : Prototile := ⟨[(0, 0), (0, 1), (1, 0)], by decide⟩ /-- A protoset containing just the L-tromino -/ def lTrominoSet : Protoset Unit := ⟨fun \_ => lTromino⟩ /- \#\# Utility Lemmas -/ /-- Each L-tromino covers exactly 3 cells -/ theorem lTromino...
\end{definition}

\begin{theorem}
  \label{tiling_5x9_valid}
  \lean{tiling_5x9_valid}
  \uses{Valid, rectangle, tiling_5x9}
  \leanok
  \proven
  \((tiling\_5x9\_valid)\)
\end{theorem}

\begin{theorem}
  \label{tileable_5x9}
  \lean{tileable_5x9}
  \uses{LTileable, rectangle, tiling_5x9, tiling_5x9_valid}
  \leanok
  \proven
  \((tileable\_5x9)\)
\end{theorem}

\begin{theorem}
  \label{tileable_9x5}
  \lean{tileable_9x5}
  \uses{LTileable, LTileable_swap_rectangle, rectangle, tileable_5x9}
  \leanok
  \proven
  The L-tromino shape, anchored at origin -/ def lTromino : Prototile := ⟨[(0, 0), (0, 1), (1, 0)], by decide⟩ /-- A protoset containing just the L-tromino -/ def lTrominoSet : Protoset Unit := ⟨fun \_ => lTromino⟩ /- \#\# Utility Lemmas -/ /-- Each L-tromino covers exactly 3 cells -/ theorem lTromino...
\end{theorem}

\begin{theorem}
  \label{tileable_6x3}
  \lean{tileable_6x3}
  \uses{LTileable, LTileable_horizontal_union, rectangle, tileable_2x3}
  \leanok
  \proven
  The L-tromino shape, anchored at origin -/ def lTromino : Prototile := ⟨[(0, 0), (0, 1), (1, 0)], by decide⟩ /-- A protoset containing just the L-tromino -/ def lTrominoSet : Protoset Unit := ⟨fun \_ => lTromino⟩ /- \#\# Utility Lemmas -/ /-- Each L-tromino covers exactly 3 cells -/ theorem lTromino...
\end{theorem}

\begin{theorem}
  \label{tileable_3x_even}
  \lean{tileable_3x_even}
  \uses{LTileable, LTileable_vertical_union, rectangle, tileable_3x2}
  \leanok
  \proven
  The L-tromino shape, anchored at origin -/ def lTromino : Prototile := ⟨[(0, 0), (0, 1), (1, 0)], by decide⟩ /-- A protoset containing just the L-tromino -/ def lTrominoSet : Protoset Unit := ⟨fun \_ => lTromino⟩ /- \#\# Utility Lemmas -/ /-- Each L-tromino covers exactly 3 cells -/ theorem lTromino...
\end{theorem}

\begin{theorem}
  \label{tileable_even_x3}
  \lean{tileable_even_x3}
  \uses{LTileable, LTileable_swap_rectangle, rectangle, tileable_3x_even}
  \leanok
  \proven
  The L-tromino shape, anchored at origin -/ def lTromino : Prototile := ⟨[(0, 0), (0, 1), (1, 0)], by decide⟩ /-- A protoset containing just the L-tromino -/ def lTrominoSet : Protoset Unit := ⟨fun \_ => lTromino⟩ /- \#\# Utility Lemmas -/ /-- Each L-tromino covers exactly 3 cells -/ theorem lTromino...
\end{theorem}

\begin{theorem}
  \label{tileable_2x_mult3}
  \lean{tileable_2x_mult3}
  \uses{LTileable, LTileable_vertical_union, empty_LTileable, rectangle,
        rectangle_n_0, tileable_2x3}
  \leanok
  \proven
  The L-tromino shape, anchored at origin -/ def lTromino : Prototile := ⟨[(0, 0), (0, 1), (1, 0)], by decide⟩ /-- A protoset containing just the L-tromino -/ def lTrominoSet : Protoset Unit := ⟨fun \_ => lTromino⟩ /- \#\# Utility Lemmas -/ /-- Each L-tromino covers exactly 3 cells -/ theorem lTromino...
\end{theorem}

\begin{theorem}
  \label{tileable_mult3_x2}
  \lean{tileable_mult3_x2}
  \uses{LTileable, LTileable_swap_rectangle, rectangle, tileable_2x_mult3}
  \leanok
  \proven
  The L-tromino shape, anchored at origin -/ def lTromino : Prototile := ⟨[(0, 0), (0, 1), (1, 0)], by decide⟩ /-- A protoset containing just the L-tromino -/ def lTrominoSet : Protoset Unit := ⟨fun \_ => lTromino⟩ /- \#\# Utility Lemmas -/ /-- Each L-tromino covers exactly 3 cells -/ theorem lTromino...
\end{theorem}

\begin{theorem}
  \label{tileable_even_mult3}
  \lean{tileable_even_mult3}
  \uses{LTileable, LTileable_horizontal_union, rectangle, tileable_2x_mult3}
  \leanok
  \proven
  The L-tromino shape, anchored at origin -/ def lTromino : Prototile := ⟨[(0, 0), (0, 1), (1, 0)], by decide⟩ /-- A protoset containing just the L-tromino -/ def lTrominoSet : Protoset Unit := ⟨fun \_ => lTromino⟩ /- \#\# Utility Lemmas -/ /-- Each L-tromino covers exactly 3 cells -/ theorem lTromino...
\end{theorem}

\begin{theorem}
  \label{tileable_3x6j}
  \lean{tileable_3x6j}
  \uses{LTileable, LTileable_vertical_union, empty_LTileable, rectangle,
        rectangle_n_0, tileable_3x6}
  \leanok
  \proven
  The L-tromino shape, anchored at origin -/ def lTromino : Prototile := ⟨[(0, 0), (0, 1), (1, 0)], by decide⟩ /-- A protoset containing just the L-tromino -/ def lTrominoSet : Protoset Unit := ⟨fun \_ => lTromino⟩ /- \#\# Utility Lemmas -/ /-- Each L-tromino covers exactly 3 cells -/ theorem lTromino...
\end{theorem}

\begin{theorem}
  \label{tileable_5x6j}
  \lean{tileable_5x6j}
  \uses{LTileable, LTileable_vertical_union, empty_LTileable, rectangle,
        rectangle_n_0, tileable_5x6}
  \leanok
  \proven
  The L-tromino shape, anchored at origin -/ def lTromino : Prototile := ⟨[(0, 0), (0, 1), (1, 0)], by decide⟩ /-- A protoset containing just the L-tromino -/ def lTrominoSet : Protoset Unit := ⟨fun \_ => lTromino⟩ /- \#\# Utility Lemmas -/ /-- Each L-tromino covers exactly 3 cells -/ theorem lTromino...
\end{theorem}

\begin{theorem}
  \label{tileable_5x_9plus6j}
  \lean{tileable_5x_9plus6j}
  \uses{LTileable, LTileable_vertical_union, rectangle, tileable_5x6j,
        tileable_5x9}
  \leanok
  \proven
  The L-tromino shape, anchored at origin -/ def lTromino : Prototile := ⟨[(0, 0), (0, 1), (1, 0)], by decide⟩ /-- A protoset containing just the L-tromino -/ def lTrominoSet : Protoset Unit := ⟨fun \_ => lTromino⟩ /- \#\# Utility Lemmas -/ /-- Each L-tromino covers exactly 3 cells -/ theorem lTromino...
\end{theorem}

\begin{theorem}
  \label{tileable_odd_x_6j}
  \lean{tileable_odd_x_6j}
  \uses{LTileable, LTileable_horizontal_union, rectangle, tileable_2x_mult3,
        tileable_3x6j, tileable_5x6j}
  \leanok
  \proven
  The L-tromino shape, anchored at origin -/ def lTromino : Prototile := ⟨[(0, 0), (0, 1), (1, 0)], by decide⟩ /-- A protoset containing just the L-tromino -/ def lTrominoSet : Protoset Unit := ⟨fun \_ => lTromino⟩ /- \#\# Utility Lemmas -/ /-- Each L-tromino covers exactly 3 cells -/ theorem lTromino...
\end{theorem}

\begin{theorem}
  \label{tileable_5x_6iplus3}
  \lean{tileable_5x_6iplus3}
  \uses{LTileable, rectangle, tileable_5x_9plus6j}
  \leanok
  \proven
  The L-tromino shape, anchored at origin -/ def lTromino : Prototile := ⟨[(0, 0), (0, 1), (1, 0)], by decide⟩ /-- A protoset containing just the L-tromino -/ def lTrominoSet : Protoset Unit := ⟨fun \_ => lTromino⟩ /- \#\# Utility Lemmas -/ /-- Each L-tromino covers exactly 3 cells -/ theorem lTromino...
\end{theorem}

\begin{theorem}
  \label{tileable_odd_ge5_x_6iplus3}
  \lean{tileable_odd_ge5_x_6iplus3}
  \uses{LTileable, LTileable_horizontal_union, rectangle, tileable_2x_mult3,
        tileable_5x_6iplus3}
  \leanok
  \proven
  The L-tromino shape, anchored at origin -/ def lTromino : Prototile := ⟨[(0, 0), (0, 1), (1, 0)], by decide⟩ /-- A protoset containing just the L-tromino -/ def lTrominoSet : Protoset Unit := ⟨fun \_ => lTromino⟩ /- \#\# Utility Lemmas -/ /-- Each L-tromino covers exactly 3 cells -/ theorem lTromino...
\end{theorem}

\begin{theorem}
  \label{tileable_odd_x_mult3}
  \lean{tileable_odd_x_mult3}
  \uses{LTileable, rectangle, tileable_odd_ge5_x_6iplus3, tileable_odd_x_6j}
  \leanok
  \proven
  The L-tromino shape, anchored at origin -/ def lTromino : Prototile := ⟨[(0, 0), (0, 1), (1, 0)], by decide⟩ /-- A protoset containing just the L-tromino -/ def lTrominoSet : Protoset Unit := ⟨fun \_ => lTromino⟩ /- \#\# Utility Lemmas -/ /-- Each L-tromino covers exactly 3 cells -/ theorem lTromino...
\end{theorem}

\begin{theorem}
  \label{rect_tileable_of_conditions}
  \lean{rect_tileable_of_conditions}
  \uses{LTileable, LTileable_swap_rectangle, RectTileableConditions, empty_LTileable,
        rectangle, rectangle_0_m, rectangle_n_0, tileable_even_mult3,
        tileable_odd_x_mult3}
  \leanok
  \proven
  The L-tromino shape, anchored at origin -/ def lTromino : Prototile := ⟨[(0, 0), (0, 1), (1, 0)], by decide⟩ /-- A protoset containing just the L-tromino -/ def lTrominoSet : Protoset Unit := ⟨fun \_ => lTromino⟩ /- \#\# Utility Lemmas -/ /-- Each L-tromino covers exactly 3 cells -/ theorem lTromino...
\end{theorem}

\begin{theorem}
  \label{conditions_of_rect_tileable}
  \lean{conditions_of_rect_tileable}
  \uses{LTileable, RectTileableConditions, not_tileable_1_by_n, not_tileable_3_by_odd,
        not_tileable_n_by_1, not_tileable_odd_by_3, rectangle}
  \leanok
  \proven
  The L-tromino shape, anchored at origin -/ def lTromino : Prototile := ⟨[(0, 0), (0, 1), (1, 0)], by decide⟩ /-- A protoset containing just the L-tromino -/ def lTrominoSet : Protoset Unit := ⟨fun \_ => lTromino⟩ /- \#\# Utility Lemmas -/ /-- Each L-tromino covers exactly 3 cells -/ theorem lTromino...
\end{theorem}

\begin{theorem}
  \label{rect_tileable_iff}
  \lean{rect_tileable_iff}
  \uses{LTileable, RectTileableConditions, conditions_of_rect_tileable, rect_tileable_of_conditions,
        rectangle}
  \leanok
  \proven
  The L-tromino shape, anchored at origin -/ def lTromino : Prototile := ⟨[(0, 0), (0, 1), (1, 0)], by decide⟩ /-- A protoset containing just the L-tromino -/ def lTrominoSet : Protoset Unit := ⟨fun \_ => lTromino⟩ /- \#\# Utility Lemmas -/ /-- Each L-tromino covers exactly 3 cells -/ theorem lTromino...
\end{theorem}

\section{Deficient Rectangles (Three-Cornered)}

\begin{definition}
  \label{cornerTR}
  \lean{cornerTR}
  \uses{Cell, rectangle}
  \leanok
  \proven
  The L-tromino shape, anchored at origin -/ def lTromino : Prototile := ⟨[(0, 0), (0, 1), (1, 0)], by decide⟩ /-- A protoset containing just the L-tromino -/ def lTrominoSet : Protoset Unit := ⟨fun \_ => lTromino⟩ /- \#\# Utility Lemmas -/ /-- Each L-tromino covers exactly 3 cells -/ theorem lTromino...
\end{definition}

\begin{definition}
  \label{cornerTR2}
  \lean{cornerTR2}
  \uses{Cell, rectangle}
  \leanok
  \proven
  The L-tromino shape, anchored at origin -/ def lTromino : Prototile := ⟨[(0, 0), (0, 1), (1, 0)], by decide⟩ /-- A protoset containing just the L-tromino -/ def lTrominoSet : Protoset Unit := ⟨fun \_ => lTromino⟩ /- \#\# Utility Lemmas -/ /-- Each L-tromino covers exactly 3 cells -/ theorem lTromino...
\end{definition}

\begin{definition}
  \label{rectangleMinusCorner}
  \lean{rectangleMinusCorner}
  \uses{Region, cornerTR, rectangle}
  \leanok
  \proven
  The L-tromino shape, anchored at origin -/ def lTromino : Prototile := ⟨[(0, 0), (0, 1), (1, 0)], by decide⟩ /-- A protoset containing just the L-tromino -/ def lTrominoSet : Protoset Unit := ⟨fun \_ => lTromino⟩ /- \#\# Utility Lemmas -/ /-- Each L-tromino covers exactly 3 cells -/ theorem lTromino...
\end{definition}

\begin{theorem}
  \label{cornerTR_mem_rectangle}
  \lean{cornerTR_mem_rectangle}
  \uses{cornerTR, cornerTR2, mem_rectangle, rectangle}
  \leanok
  \proven
  The L-tromino shape, anchored at origin -/ def lTromino : Prototile := ⟨[(0, 0), (0, 1), (1, 0)], by decide⟩ /-- A protoset containing just the L-tromino -/ def lTrominoSet : Protoset Unit := ⟨fun \_ => lTromino⟩ /- \#\# Utility Lemmas -/ /-- Each L-tromino covers exactly 3 cells -/ theorem lTromino...
\end{theorem}

\begin{theorem}
  \label{cornerTR2_mem_rectangle}
  \lean{cornerTR2_mem_rectangle}
  \uses{cornerTR2, mem_rectangle, rectangle}
  \leanok
  \proven
  The L-tromino shape, anchored at origin -/ def lTromino : Prototile := ⟨[(0, 0), (0, 1), (1, 0)], by decide⟩ /-- A protoset containing just the L-tromino -/ def lTrominoSet : Protoset Unit := ⟨fun \_ => lTromino⟩ /- \#\# Utility Lemmas -/ /-- Each L-tromino covers exactly 3 cells -/ theorem lTromino...
\end{theorem}

\begin{theorem}
  \label{rectangleMinusCorner_card}
  \lean{rectangleMinusCorner_card}
  \uses{cornerTR, cornerTR_mem_rectangle, rectangle, rectangleMinusCorner}
  \leanok
  \proven
  The L-tromino shape, anchored at origin -/ def lTromino : Prototile := ⟨[(0, 0), (0, 1), (1, 0)], by decide⟩ /-- A protoset containing just the L-tromino -/ def lTrominoSet : Protoset Unit := ⟨fun \_ => lTromino⟩ /- \#\# Utility Lemmas -/ /-- Each L-tromino covers exactly 3 cells -/ theorem lTromino...
\end{theorem}

\begin{definition}
  \label{rectangleMinusCorner_rexp}
  \lean{rectangleMinusCorner_rexp}
  \uses{RExp_r, rect, rectangle}
  \leanok
  \proven
  The L-tromino shape, anchored at origin -/ def lTromino : Prototile := ⟨[(0, 0), (0, 1), (1, 0)], by decide⟩ /-- A protoset containing just the L-tromino -/ def lTrominoSet : Protoset Unit := ⟨fun \_ => lTromino⟩ /- \#\# Utility Lemmas -/ /-- Each L-tromino covers exactly 3 cells -/ theorem lTromino...
\end{definition}

\begin{theorem}
  \label{rectangle_eq_rect_coe}
  \lean{rectangle_eq_rect_coe}
  \uses{Cell, mem_rectangle, rect, rectangle,
        rectangleMinusCorner, rectangleMinusCorner_rexp}
  \leanok
  \proven
  The L-tromino shape, anchored at origin -/ def lTromino : Prototile := ⟨[(0, 0), (0, 1), (1, 0)], by decide⟩ /-- A protoset containing just the L-tromino -/ def lTrominoSet : Protoset Unit := ⟨fun \_ => lTromino⟩ /- \#\# Utility Lemmas -/ /-- Each L-tromino covers exactly 3 cells -/ theorem lTromino...
\end{theorem}

\begin{theorem}
  \label{rectangleMinusCorner_rexp_eval}
  \lean{rectangleMinusCorner_rexp_eval}
  \uses{Cell, RExp_eval, cornerTR, mem_rectangle,
        rectangle, rectangleMinusCorner, rectangleMinusCorner_rexp, rectangle_eq_rect_coe}
  \leanok
  \proven
  The L-tromino shape, anchored at origin -/ def lTromino : Prototile := ⟨[(0, 0), (0, 1), (1, 0)], by decide⟩ /-- A protoset containing just the L-tromino -/ def lTrominoSet : Protoset Unit := ⟨fun \_ => lTromino⟩ /- \#\# Utility Lemmas -/ /-- Each L-tromino covers exactly 3 cells -/ theorem lTromino...
\end{theorem}

\begin{theorem}
  \label{swap_rectangleMinusCorner_rexp}
  \lean{swap_rectangleMinusCorner_rexp}
  \uses{RExp_eval, rectangle, rectangleMinusCorner_rexp}
  \leanok
  \proven
  The L-tromino shape, anchored at origin -/ def lTromino : Prototile := ⟨[(0, 0), (0, 1), (1, 0)], by decide⟩ /-- A protoset containing just the L-tromino -/ def lTrominoSet : Protoset Unit := ⟨fun \_ => lTromino⟩ /- \#\# Utility Lemmas -/ /-- Each L-tromino covers exactly 3 cells -/ theorem lTromino...
\end{theorem}

\begin{theorem}
  \label{rectangleMinusCorner_split_horizontal_rexp}
  \lean{rectangleMinusCorner_split_horizontal_rexp}
  \uses{RExp_eval, RExp_r, rectangle, rectangleMinusCorner_rexp}
  \leanok
  \proven
  The L-tromino shape, anchored at origin -/ def lTromino : Prototile := ⟨[(0, 0), (0, 1), (1, 0)], by decide⟩ /-- A protoset containing just the L-tromino -/ def lTrominoSet : Protoset Unit := ⟨fun \_ => lTromino⟩ /- \#\# Utility Lemmas -/ /-- Each L-tromino covers exactly 3 cells -/ theorem lTromino...
\end{theorem}

\begin{theorem}
  \label{rectangleMinusCorner_split_vertical_rexp}
  \lean{rectangleMinusCorner_split_vertical_rexp}
  \uses{RExp_eval, RExp_r, rectangleMinusCorner_rexp, reflect,
        swapRegion}
  \leanok
  \proven
  The L-tromino shape, anchored at origin -/ def lTromino : Prototile := ⟨[(0, 0), (0, 1), (1, 0)], by decide⟩ /-- A protoset containing just the L-tromino -/ def lTrominoSet : Protoset Unit := ⟨fun \_ => lTromino⟩ /- \#\# Utility Lemmas -/ /-- Each L-tromino covers exactly 3 cells -/ theorem lTromino...
\end{theorem}

\begin{theorem}
  \label{swapRegion_eq_reflect_coe}
  \lean{swapRegion_eq_reflect_coe}
  \uses{Cell, Region, reflect, swapCell,
        swapCell_involutive, swapRegion, translate, translateRegion}
  \leanok
  \proven
  The L-tromino shape, anchored at origin -/ def lTromino : Prototile := ⟨[(0, 0), (0, 1), (1, 0)], by decide⟩ /-- A protoset containing just the L-tromino -/ def lTrominoSet : Protoset Unit := ⟨fun \_ => lTromino⟩ /- \#\# Utility Lemmas -/ /-- Each L-tromino covers exactly 3 cells -/ theorem lTromino...
\end{theorem}

\begin{theorem}
  \label{translateRegion_eq_translate_coe}
  \lean{translateRegion_eq_translate_coe}
  \uses{Cell, Region, rectangle, translate,
        translateCell, translateRegion}
  \leanok
  \proven
  The L-tromino shape, anchored at origin -/ def lTromino : Prototile := ⟨[(0, 0), (0, 1), (1, 0)], by decide⟩ /-- A protoset containing just the L-tromino -/ def lTrominoSet : Protoset Unit := ⟨fun \_ => lTromino⟩ /- \#\# Utility Lemmas -/ /-- Each L-tromino covers exactly 3 cells -/ theorem lTromino...
\end{theorem}

\begin{theorem}
  \label{swap_rectangleMinusCorner}
  \lean{swap_rectangleMinusCorner}
  \uses{Cell, RExp_eval, rectangle, rectangleMinusCorner,
        rectangleMinusCorner_rexp, rectangleMinusCorner_rexp_eval, reflect, swapRegion,
        swapRegion_eq_reflect_coe, swap_rectangleMinusCorner_rexp}
  \leanok
  \proven
  The L-tromino shape, anchored at origin -/ def lTromino : Prototile := ⟨[(0, 0), (0, 1), (1, 0)], by decide⟩ /-- A protoset containing just the L-tromino -/ def lTrominoSet : Protoset Unit := ⟨fun \_ => lTromino⟩ /- \#\# Utility Lemmas -/ /-- Each L-tromino covers exactly 3 cells -/ theorem lTromino...
\end{theorem}

\begin{definition}
  \label{rectangleMinus2Corner}
  \lean{rectangleMinus2Corner}
  \uses{Region, cornerTR, cornerTR2, rectangle}
  \leanok
  \proven
  The L-tromino shape, anchored at origin -/ def lTromino : Prototile := ⟨[(0, 0), (0, 1), (1, 0)], by decide⟩ /-- A protoset containing just the L-tromino -/ def lTrominoSet : Protoset Unit := ⟨fun \_ => lTromino⟩ /- \#\# Utility Lemmas -/ /-- Each L-tromino covers exactly 3 cells -/ theorem lTromino...
\end{definition}

\begin{theorem}
  \label{rectangleMinus2Corner_card}
  \lean{rectangleMinus2Corner_card}
  \uses{cornerTR, cornerTR2, cornerTR2_mem_rectangle, cornerTR_mem_rectangle,
        rectangle, rectangleMinus2Corner}
  \leanok
  \proven
  The L-tromino shape, anchored at origin -/ def lTromino : Prototile := ⟨[(0, 0), (0, 1), (1, 0)], by decide⟩ /-- A protoset containing just the L-tromino -/ def lTrominoSet : Protoset Unit := ⟨fun \_ => lTromino⟩ /- \#\# Utility Lemmas -/ /-- Each L-tromino covers exactly 3 cells -/ theorem lTromino...
\end{theorem}

\begin{theorem}
  \label{swapCell_cornerTR}
  \lean{swapCell_cornerTR}
  \uses{cornerTR, swapCell}
  \leanok
  \proven
  The L-tromino shape, anchored at origin -/ def lTromino : Prototile := ⟨[(0, 0), (0, 1), (1, 0)], by decide⟩ /-- A protoset containing just the L-tromino -/ def lTrominoSet : Protoset Unit := ⟨fun \_ => lTromino⟩ /- \#\# Utility Lemmas -/ /-- Each L-tromino covers exactly 3 cells -/ theorem lTromino...
\end{theorem}

\begin{theorem}
  \label{LTileable_swap_rectangleMinusCorner}
  \lean{LTileable_swap_rectangleMinusCorner}
  \uses{LTileable, LTileable_swap, rectangle, rectangleMinusCorner,
        swap_rectangleMinusCorner}
  \leanok
  \proven
  The L-tromino shape, anchored at origin -/ def lTromino : Prototile := ⟨[(0, 0), (0, 1), (1, 0)], by decide⟩ /-- A protoset containing just the L-tromino -/ def lTrominoSet : Protoset Unit := ⟨fun \_ => lTromino⟩ /- \#\# Utility Lemmas -/ /-- Each L-tromino covers exactly 3 cells -/ theorem lTromino...
\end{theorem}

\section{Splitting three-cornered rectangles}

\begin{theorem}
  \label{rectangleMinusCorner_split_horizontal}
  \lean{rectangleMinusCorner_split_horizontal}
  \uses{Cell, RExp_eval, RExp_r, rectangle,
        rectangleMinusCorner, rectangleMinusCorner_rexp, rectangleMinusCorner_rexp_eval, rectangleMinusCorner_split_horizontal_rexp,
        rectangle_eq_rect_coe, translateRegion, translateRegion_eq_translate_coe}
  \leanok
  \proven
  The L-tromino shape, anchored at origin -/ def lTromino : Prototile := ⟨[(0, 0), (0, 1), (1, 0)], by decide⟩ /-- A protoset containing just the L-tromino -/ def lTrominoSet : Protoset Unit := ⟨fun \_ => lTromino⟩ /- \#\# Utility Lemmas -/ /-- Each L-tromino covers exactly 3 cells -/ theorem lTromino...
\end{theorem}

\begin{theorem}
  \label{rectangleMinusCorner_split_vertical}
  \lean{rectangleMinusCorner_split_vertical}
  \uses{Cell, RExp_eval, RExp_r, rectangle,
        rectangleMinusCorner, rectangleMinusCorner_rexp, rectangleMinusCorner_rexp_eval, rectangleMinusCorner_split_vertical_rexp,
        rectangle_eq_rect_coe, translateRegion, translateRegion_eq_translate_coe}
  \leanok
  \proven
  The L-tromino shape, anchored at origin -/ def lTromino : Prototile := ⟨[(0, 0), (0, 1), (1, 0)], by decide⟩ /-- A protoset containing just the L-tromino -/ def lTrominoSet : Protoset Unit := ⟨fun \_ => lTromino⟩ /- \#\# Utility Lemmas -/ /-- Each L-tromino covers exactly 3 cells -/ theorem lTromino...
\end{theorem}

\section{Simple mod-3 helper lemmas}

\begin{theorem}
  \label{mod3_ne_zero_of_mul_mod3_ne_zero_left}
  \lean{mod3_ne_zero_of_mul_mod3_ne_zero_left}
  \leanok
  \proven
  The L-tromino shape, anchored at origin -/ def lTromino : Prototile := ⟨[(0, 0), (0, 1), (1, 0)], by decide⟩ /-- A protoset containing just the L-tromino -/ def lTrominoSet : Protoset Unit := ⟨fun \_ => lTromino⟩ /- \#\# Utility Lemmas -/ /-- Each L-tromino covers exactly 3 cells -/ theorem lTromino...
\end{theorem}

\begin{theorem}
  \label{mod3_eq_one_or_two_of_ne_zero}
  \lean{mod3_eq_one_or_two_of_ne_zero}
  \leanok
  \proven
  The L-tromino shape, anchored at origin -/ def lTromino : Prototile := ⟨[(0, 0), (0, 1), (1, 0)], by decide⟩ /-- A protoset containing just the L-tromino -/ def lTrominoSet : Protoset Unit := ⟨fun \_ => lTromino⟩ /- \#\# Utility Lemmas -/ /-- Each L-tromino covers exactly 3 cells -/ theorem lTromino...
\end{theorem}

\begin{theorem}
  \label{mod3_both_one_or_two_of_area_mod3_zero}
  \lean{mod3_both_one_or_two_of_area_mod3_zero}
  \uses{mod3_eq_one_or_two_of_ne_zero, mod3_ne_zero_of_mul_mod3_ne_zero_left}
  \leanok
  \proven
  The L-tromino shape, anchored at origin -/ def lTromino : Prototile := ⟨[(0, 0), (0, 1), (1, 0)], by decide⟩ /-- A protoset containing just the L-tromino -/ def lTrominoSet : Protoset Unit := ⟨fun \_ => lTromino⟩ /- \#\# Utility Lemmas -/ /-- Each L-tromino covers exactly 3 cells -/ theorem lTromino...
\end{theorem}

\begin{theorem}
  \label{mod3_one_two_or_two_one_of_area_mod3_two}
  \lean{mod3_one_two_or_two_one_of_area_mod3_two}
  \uses{mod3_eq_one_or_two_of_ne_zero, mod3_ne_zero_of_mul_mod3_ne_zero_left, rectangle}
  \leanok
  \proven
  The L-tromino shape, anchored at origin -/ def lTromino : Prototile := ⟨[(0, 0), (0, 1), (1, 0)], by decide⟩ /-- A protoset containing just the L-tromino -/ def lTrominoSet : Protoset Unit := ⟨fun \_ => lTromino⟩ /- \#\# Utility Lemmas -/ /-- Each L-tromino covers exactly 3 cells -/ theorem lTromino...
\end{theorem}

\begin{theorem}
  \label{rectMinusCorner_tileable_area_div_3}
  \lean{rectMinusCorner_tileable_area_div_3}
  \uses{LTileable, rectangle, rectangleMinusCorner, rectangleMinusCorner_card}
  \leanok
  \proven
  The L-tromino shape, anchored at origin -/ def lTromino : Prototile := ⟨[(0, 0), (0, 1), (1, 0)], by decide⟩ /-- A protoset containing just the L-tromino -/ def lTrominoSet : Protoset Unit := ⟨fun \_ => lTromino⟩ /- \#\# Utility Lemmas -/ /-- Each L-tromino covers exactly 3 cells -/ theorem lTromino...
\end{theorem}

\begin{definition}
  \label{tiling_2x2_minus}
  \lean{tiling_2x2_minus}
  \uses{TileSet, lTrominoSet, mkPlacedTile, mkTileSet}
  \leanok
  \proven
  The L-tromino shape, anchored at origin -/ def lTromino : Prototile := ⟨[(0, 0), (0, 1), (1, 0)], by decide⟩ /-- A protoset containing just the L-tromino -/ def lTrominoSet : Protoset Unit := ⟨fun \_ => lTromino⟩ /- \#\# Utility Lemmas -/ /-- Each L-tromino covers exactly 3 cells -/ theorem lTromino...
\end{definition}

\begin{theorem}
  \label{tiling_2x2_minus_valid}
  \lean{tiling_2x2_minus_valid}
  \uses{Valid, rectangle, rectangleMinusCorner, tiling_2x2_minus}
  \leanok
  \proven
  \((tiling\_2x2\_minus\_valid)\)
\end{theorem}

\begin{theorem}
  \label{tileable_2x2_minus}
  \lean{tileable_2x2_minus}
  \uses{LTileable, rectangle, rectangleMinusCorner, tiling_2x2_minus,
        tiling_2x2_minus_valid}
  \leanok
  \proven
  The L-tromino shape, anchored at origin -/ def lTromino : Prototile := ⟨[(0, 0), (0, 1), (1, 0)], by decide⟩ /-- A protoset containing just the L-tromino -/ def lTrominoSet : Protoset Unit := ⟨fun \_ => lTromino⟩ /- \#\# Utility Lemmas -/ /-- Each L-tromino covers exactly 3 cells -/ theorem lTromino...
\end{theorem}

\begin{definition}
  \label{tiling_5x2_minus}
  \lean{tiling_5x2_minus}
  \uses{TileSet, lTrominoSet}
  \leanok
  \proven
  The L-tromino shape, anchored at origin -/ def lTromino : Prototile := ⟨[(0, 0), (0, 1), (1, 0)], by decide⟩ /-- A protoset containing just the L-tromino -/ def lTrominoSet : Protoset Unit := ⟨fun \_ => lTromino⟩ /- \#\# Utility Lemmas -/ /-- Each L-tromino covers exactly 3 cells -/ theorem lTromino...
\end{definition}

\begin{theorem}
  \label{tiling_5x2_minus_valid}
  \lean{tiling_5x2_minus_valid}
  \uses{Valid, rectangle, rectangleMinusCorner, tiling_5x2_minus}
  \leanok
  \proven
  \((tiling\_5x2\_minus\_valid)\)
\end{theorem}

\begin{theorem}
  \label{tileable_5x2_minus}
  \lean{tileable_5x2_minus}
  \uses{LTileable, rectangle, rectangleMinusCorner, tiling_5x2_minus,
        tiling_5x2_minus_valid}
  \leanok
  \proven
  The L-tromino shape, anchored at origin -/ def lTromino : Prototile := ⟨[(0, 0), (0, 1), (1, 0)], by decide⟩ /-- A protoset containing just the L-tromino -/ def lTrominoSet : Protoset Unit := ⟨fun \_ => lTromino⟩ /- \#\# Utility Lemmas -/ /-- Each L-tromino covers exactly 3 cells -/ theorem lTromino...
\end{theorem}

\begin{definition}
  \label{tiling_4x4_minus}
  \lean{tiling_4x4_minus}
  \uses{TileSet, lTrominoSet}
  \leanok
  \proven
  The L-tromino shape, anchored at origin -/ def lTromino : Prototile := ⟨[(0, 0), (0, 1), (1, 0)], by decide⟩ /-- A protoset containing just the L-tromino -/ def lTrominoSet : Protoset Unit := ⟨fun \_ => lTromino⟩ /- \#\# Utility Lemmas -/ /-- Each L-tromino covers exactly 3 cells -/ theorem lTromino...
\end{definition}

\begin{theorem}
  \label{tiling_4x4_minus_valid}
  \lean{tiling_4x4_minus_valid}
  \uses{Valid, rectangle, rectangleMinusCorner, tiling_4x4_minus}
  \leanok
  \proven
  \((tiling\_4x4\_minus\_valid)\)
\end{theorem}

\begin{theorem}
  \label{tileable_4x4_minus}
  \lean{tileable_4x4_minus}
  \uses{LTileable, rectangle, rectangleMinusCorner, tiling_2x3,
        tiling_4x4_minus, tiling_4x4_minus_valid}
  \leanok
  \proven
  The L-tromino shape, anchored at origin -/ def lTromino : Prototile := ⟨[(0, 0), (0, 1), (1, 0)], by decide⟩ /-- A protoset containing just the L-tromino -/ def lTrominoSet : Protoset Unit := ⟨fun \_ => lTromino⟩ /- \#\# Utility Lemmas -/ /-- Each L-tromino covers exactly 3 cells -/ theorem lTromino...
\end{theorem}

\begin{definition}
  \label{tiling_5x5_minus}
  \lean{tiling_5x5_minus}
  \uses{TileSet, lTrominoSet, rectangle, tiling_2x3}
  \leanok
  \proven
  The L-tromino shape, anchored at origin -/ def lTromino : Prototile := ⟨[(0, 0), (0, 1), (1, 0)], by decide⟩ /-- A protoset containing just the L-tromino -/ def lTrominoSet : Protoset Unit := ⟨fun \_ => lTromino⟩ /- \#\# Utility Lemmas -/ /-- Each L-tromino covers exactly 3 cells -/ theorem lTromino...
\end{definition}

\begin{theorem}
  \label{tiling_5x5_minus_valid}
  \lean{tiling_5x5_minus_valid}
  \uses{Valid, rectangle, rectangleMinusCorner, tiling_5x5_minus}
  \leanok
  \proven
  \((tiling\_5x5\_minus\_valid)\)
\end{theorem}

\begin{theorem}
  \label{tileable_5x5_minus}
  \lean{tileable_5x5_minus}
  \uses{LTileable, rectangle, rectangleMinusCorner, tiling_5x5_minus,
        tiling_5x5_minus_valid}
  \leanok
  \proven
  The L-tromino shape, anchored at origin -/ def lTromino : Prototile := ⟨[(0, 0), (0, 1), (1, 0)], by decide⟩ /-- A protoset containing just the L-tromino -/ def lTrominoSet : Protoset Unit := ⟨fun \_ => lTromino⟩ /- \#\# Utility Lemmas -/ /-- Each L-tromino covers exactly 3 cells -/ theorem lTromino...
\end{theorem}

\begin{definition}
  \label{tiling_7x7_minus}
  \lean{tiling_7x7_minus}
  \uses{TileSet, lTrominoSet}
  \leanok
  \proven
  The L-tromino shape, anchored at origin -/ def lTromino : Prototile := ⟨[(0, 0), (0, 1), (1, 0)], by decide⟩ /-- A protoset containing just the L-tromino -/ def lTrominoSet : Protoset Unit := ⟨fun \_ => lTromino⟩ /- \#\# Utility Lemmas -/ /-- Each L-tromino covers exactly 3 cells -/ theorem lTromino...
\end{definition}

\begin{theorem}
  \label{tiling_7x7_minus_valid}
  \lean{tiling_7x7_minus_valid}
  \uses{Valid, rectangle, rectangleMinusCorner, tiling_7x7_minus}
  \leanok
  \proven
  \((tiling\_7x7\_minus\_valid)\)
\end{theorem}

\begin{theorem}
  \label{tileable_7x7_minus}
  \lean{tileable_7x7_minus}
  \uses{LTileable, rectangle, rectangleMinusCorner, tiling_7x7_minus,
        tiling_7x7_minus_valid}
  \leanok
  \proven
  The L-tromino shape, anchored at origin -/ def lTromino : Prototile := ⟨[(0, 0), (0, 1), (1, 0)], by decide⟩ /-- A protoset containing just the L-tromino -/ def lTrominoSet : Protoset Unit := ⟨fun \_ => lTromino⟩ /- \#\# Utility Lemmas -/ /-- Each L-tromino covers exactly 3 cells -/ theorem lTromino...
\end{theorem}

\section{Combining rectangles with three-cornered rectangles}

\begin{theorem}
  \label{rectangleMinusCorner_split_vertical_disjoint}
  \lean{rectangleMinusCorner_split_vertical_disjoint}
  \uses{mem_rectangle, rectangle, rectangleMinusCorner, translateCell,
        translateRegion}
  \leanok
  \proven
  The L-tromino shape, anchored at origin -/ def lTromino : Prototile := ⟨[(0, 0), (0, 1), (1, 0)], by decide⟩ /-- A protoset containing just the L-tromino -/ def lTrominoSet : Protoset Unit := ⟨fun \_ => lTromino⟩ /- \#\# Utility Lemmas -/ /-- Each L-tromino covers exactly 3 cells -/ theorem lTromino...
\end{theorem}

\begin{theorem}
  \label{rectangleMinusCorner_split_horizontal_disjoint}
  \lean{rectangleMinusCorner_split_horizontal_disjoint}
  \uses{mem_rectangle, rectangle, rectangleMinusCorner, translateCell,
        translateRegion}
  \leanok
  \proven
  The L-tromino shape, anchored at origin -/ def lTromino : Prototile := ⟨[(0, 0), (0, 1), (1, 0)], by decide⟩ /-- A protoset containing just the L-tromino -/ def lTrominoSet : Protoset Unit := ⟨fun \_ => lTromino⟩ /- \#\# Utility Lemmas -/ /-- Each L-tromino covers exactly 3 cells -/ theorem lTromino...
\end{theorem}

\begin{theorem}
  \label{LTileable_vertical_union_rect_minus}
  \lean{LTileable_vertical_union_rect_minus}
  \uses{LTileable, LTileable_translate, Tileable_union, lTrominoSet,
        rectangle, rectangleMinusCorner, rectangleMinusCorner_split_vertical, rectangleMinusCorner_split_vertical_disjoint}
  \leanok
  \proven
  The L-tromino shape, anchored at origin -/ def lTromino : Prototile := ⟨[(0, 0), (0, 1), (1, 0)], by decide⟩ /-- A protoset containing just the L-tromino -/ def lTrominoSet : Protoset Unit := ⟨fun \_ => lTromino⟩ /- \#\# Utility Lemmas -/ /-- Each L-tromino covers exactly 3 cells -/ theorem lTromino...
\end{theorem}

\begin{theorem}
  \label{LTileable_horizontal_union_rect_minus}
  \lean{LTileable_horizontal_union_rect_minus}
  \uses{LTileable, LTileable_translate, Tileable_union, lTrominoSet,
        rectangle, rectangleMinusCorner, rectangleMinusCorner_split_horizontal, rectangleMinusCorner_split_horizontal_disjoint}
  \leanok
  \proven
  The L-tromino shape, anchored at origin -/ def lTromino : Prototile := ⟨[(0, 0), (0, 1), (1, 0)], by decide⟩ /-- A protoset containing just the L-tromino -/ def lTrominoSet : Protoset Unit := ⟨fun \_ => lTromino⟩ /- \#\# Utility Lemmas -/ /-- Each L-tromino covers exactly 3 cells -/ theorem lTromino...
\end{theorem}

\begin{theorem}
  \label{tileable_3kplus2_x2_minus}
  \lean{tileable_3kplus2_x2_minus}
  \uses{LTileable, LTileable_horizontal_union_rect_minus, rectangle, rectangleMinusCorner,
        tileable_2x2_minus, tileable_mult3_x2}
  \leanok
  \proven
  The L-tromino shape, anchored at origin -/ def lTromino : Prototile := ⟨[(0, 0), (0, 1), (1, 0)], by decide⟩ /-- A protoset containing just the L-tromino -/ def lTrominoSet : Protoset Unit := ⟨fun \_ => lTromino⟩ /- \#\# Utility Lemmas -/ /-- Each L-tromino covers exactly 3 cells -/ theorem lTromino...
\end{theorem}

\begin{theorem}
  \label{tileable_4x7_minus}
  \lean{tileable_4x7_minus}
  \uses{LTileable, LTileable_vertical_union_rect_minus, rectangle, rectangleMinusCorner,
        tileable_4x4_minus, tileable_even_x3}
  \leanok
  \proven
  The L-tromino shape, anchored at origin -/ def lTromino : Prototile := ⟨[(0, 0), (0, 1), (1, 0)], by decide⟩ /-- A protoset containing just the L-tromino -/ def lTrominoSet : Protoset Unit := ⟨fun \_ => lTromino⟩ /- \#\# Utility Lemmas -/ /-- Each L-tromino covers exactly 3 cells -/ theorem lTromino...
\end{theorem}

\begin{theorem}
  \label{tileable_5x_6kplus5_minus}
  \lean{tileable_5x_6kplus5_minus}
  \uses{LTileable, LTileable_vertical_union_rect_minus, rectangle, rectangleMinusCorner,
        tileable_5x5_minus, tileable_5x6j}
  \leanok
  \proven
  The L-tromino shape, anchored at origin -/ def lTromino : Prototile := ⟨[(0, 0), (0, 1), (1, 0)], by decide⟩ /-- A protoset containing just the L-tromino -/ def lTrominoSet : Protoset Unit := ⟨fun \_ => lTromino⟩ /- \#\# Utility Lemmas -/ /-- Each L-tromino covers exactly 3 cells -/ theorem lTromino...
\end{theorem}

\begin{theorem}
  \label{tileable_5x_6kplus2_minus}
  \lean{tileable_5x_6kplus2_minus}
  \uses{LTileable, LTileable_vertical_union_rect_minus, rectangle, rectangleMinusCorner,
        tileable_5x2_minus, tileable_5x6j}
  \leanok
  \proven
  The L-tromino shape, anchored at origin -/ def lTromino : Prototile := ⟨[(0, 0), (0, 1), (1, 0)], by decide⟩ /-- A protoset containing just the L-tromino -/ def lTrominoSet : Protoset Unit := ⟨fun \_ => lTromino⟩ /- \#\# Utility Lemmas -/ /-- Each L-tromino covers exactly 3 cells -/ theorem lTromino...
\end{theorem}

\begin{theorem}
  \label{tileable_5x_3kplus2_minus}
  \lean{tileable_5x_3kplus2_minus}
  \uses{LTileable, rectangle, rectangleMinusCorner, tileable_5x_6kplus2_minus,
        tileable_5x_6kplus5_minus}
  \leanok
  \proven
  The L-tromino shape, anchored at origin -/ def lTromino : Prototile := ⟨[(0, 0), (0, 1), (1, 0)], by decide⟩ /-- A protoset containing just the L-tromino -/ def lTrominoSet : Protoset Unit := ⟨fun \_ => lTromino⟩ /- \#\# Utility Lemmas -/ /-- Each L-tromino covers exactly 3 cells -/ theorem lTromino...
\end{theorem}

\begin{theorem}
  \label{tileable_rectMinusCorner_mod2_jk_ge2}
  \lean{tileable_rectMinusCorner_mod2_jk_ge2}
  \uses{LTileable, LTileable_vertical_union_rect_minus, RectTileableConditions, rect_tileable_iff,
        rectangle, rectangleMinusCorner, tileable_3kplus2_x2_minus}
  \leanok
  \proven
  The L-tromino shape, anchored at origin -/ def lTromino : Prototile := ⟨[(0, 0), (0, 1), (1, 0)], by decide⟩ /-- A protoset containing just the L-tromino -/ def lTrominoSet : Protoset Unit := ⟨fun \_ => lTromino⟩ /- \#\# Utility Lemmas -/ /-- Each L-tromino covers exactly 3 cells -/ theorem lTromino...
\end{theorem}

\section{Helper lemmas for the mod 3 = 1 case}

\begin{theorem}
  \label{tileable_4x_rectMinus_3kplus1}
  \lean{tileable_4x_rectMinus_3kplus1}
  \uses{LTileable, LTileable_vertical_union_rect_minus, rectangle, rectangleMinusCorner,
        tileable_4x4_minus, tileable_even_x3}
  \leanok
  \proven
  The L-tromino shape, anchored at origin -/ def lTromino : Prototile := ⟨[(0, 0), (0, 1), (1, 0)], by decide⟩ /-- A protoset containing just the L-tromino -/ def lTrominoSet : Protoset Unit := ⟨fun \_ => lTromino⟩ /- \#\# Utility Lemmas -/ /-- Each L-tromino covers exactly 3 cells -/ theorem lTromino...
\end{theorem}

\begin{theorem}
  \label{tileable_rect_mod1_by_mult3}
  \lean{tileable_rect_mod1_by_mult3}
  \uses{LTileable, rect_tileable_iff, rectangle}
  \leanok
  \proven
  The L-tromino shape, anchored at origin -/ def lTromino : Prototile := ⟨[(0, 0), (0, 1), (1, 0)], by decide⟩ /-- A protoset containing just the L-tromino -/ def lTrominoSet : Protoset Unit := ⟨fun \_ => lTromino⟩ /- \#\# Utility Lemmas -/ /-- Each L-tromino covers exactly 3 cells -/ theorem lTromino...
\end{theorem}

\begin{theorem}
  \label{tileable_rectMinusCorner_mod1_jk_ge}
  \lean{tileable_rectMinusCorner_mod1_jk_ge}
  \uses{LTileable, LTileable_swap_rectangleMinusCorner, LTileable_vertical_union_rect_minus, rectangle,
        rectangleMinusCorner, tileable_4x_rectMinus_3kplus1, tileable_rect_mod1_by_mult3}
  \leanok
  \proven
  The L-tromino shape, anchored at origin -/ def lTromino : Prototile := ⟨[(0, 0), (0, 1), (1, 0)], by decide⟩ /-- A protoset containing just the L-tromino -/ def lTrominoSet : Protoset Unit := ⟨fun \_ => lTromino⟩ /- \#\# Utility Lemmas -/ /-- Each L-tromino covers exactly 3 cells -/ theorem lTromino...
\end{theorem}

\begin{theorem}
  \label{tileable_rectMinusCorner_mod1_width4}
  \lean{tileable_rectMinusCorner_mod1_width4}
  \uses{LTileable, LTileable_horizontal_union_rect_minus, rect_tileable_iff, rectangle,
        rectangleMinusCorner, tileable_4x4_minus, translateRegion}
  \leanok
  \proven
  The L-tromino shape, anchored at origin -/ def lTromino : Prototile := ⟨[(0, 0), (0, 1), (1, 0)], by decide⟩ /-- A protoset containing just the L-tromino -/ def lTrominoSet : Protoset Unit := ⟨fun \_ => lTromino⟩ /- \#\# Utility Lemmas -/ /-- Each L-tromino covers exactly 3 cells -/ theorem lTromino...
\end{theorem}

\begin{theorem}
  \label{tileable_rectMinusCorner_mod1_recurrence_k_ge3}
  \lean{tileable_rectMinusCorner_mod1_recurrence_k_ge3}
  \uses{LTileable, LTileable_horizontal_union_rect_minus, RectTileableConditions, rect_tileable_iff,
        rectangle, rectangleMinusCorner, tileable_4x_rectMinus_3kplus1}
  \leanok
  \proven
  The L-tromino shape, anchored at origin -/ def lTromino : Prototile := ⟨[(0, 0), (0, 1), (1, 0)], by decide⟩ /-- A protoset containing just the L-tromino -/ def lTrominoSet : Protoset Unit := ⟨fun \_ => lTromino⟩ /- \#\# Utility Lemmas -/ /-- Each L-tromino covers exactly 3 cells -/ theorem lTromino...
\end{theorem}

\begin{theorem}
  \label{tileable_rectMinusCorner_mod1_case}
  \lean{tileable_rectMinusCorner_mod1_case}
  \uses{LTileable, LTileable_swap_rectangleMinusCorner, rectangleMinusCorner, tileable_4x4_minus,
        tileable_4x7_minus, tileable_7x7_minus, tileable_rectMinusCorner_mod1_jk_ge, tileable_rectMinusCorner_mod1_recurrence_k_ge3}
  \leanok
  \proven
  The L-tromino shape, anchored at origin -/ def lTromino : Prototile := ⟨[(0, 0), (0, 1), (1, 0)], by decide⟩ /-- A protoset containing just the L-tromino -/ def lTrominoSet : Protoset Unit := ⟨fun \_ => lTromino⟩ /- \#\# Utility Lemmas -/ /-- Each L-tromino covers exactly 3 cells -/ theorem lTromino...
\end{theorem}

\begin{theorem}
  \label{tileable_rectMinusCorner_mod2_case}
  \lean{tileable_rectMinusCorner_mod2_case}
  \uses{LTileable, LTileable_swap_rectangleMinusCorner, rectangle, rectangleMinusCorner,
        tileable_2x2_minus, tileable_3kplus2_x2_minus, tileable_5x2_minus, tileable_5x5_minus,
        tileable_5x_3kplus2_minus, tileable_rectMinusCorner_mod2_jk_ge2}
  \leanok
  \proven
  The L-tromino shape, anchored at origin -/ def lTromino : Prototile := ⟨[(0, 0), (0, 1), (1, 0)], by decide⟩ /-- A protoset containing just the L-tromino -/ def lTrominoSet : Protoset Unit := ⟨fun \_ => lTromino⟩ /- \#\# Utility Lemmas -/ /-- Each L-tromino covers exactly 3 cells -/ theorem lTromino...
\end{theorem}

\begin{theorem}
  \label{rectMinusCorner_tileable_of_area_m_le_n}
  \lean{rectMinusCorner_tileable_of_area_m_le_n}
  \uses{LTileable, mod3_both_one_or_two_of_area_mod3_zero, rectangle, rectangleMinusCorner,
        tileable_rectMinusCorner_mod1_case, tileable_rectMinusCorner_mod2_case}
  \leanok
  \proven
  The L-tromino shape, anchored at origin -/ def lTromino : Prototile := ⟨[(0, 0), (0, 1), (1, 0)], by decide⟩ /-- A protoset containing just the L-tromino -/ def lTrominoSet : Protoset Unit := ⟨fun \_ => lTromino⟩ /- \#\# Utility Lemmas -/ /-- Each L-tromino covers exactly 3 cells -/ theorem lTromino...
\end{theorem}

\begin{theorem}
  \label{rectMinusCorner_tileable_iff}
  \lean{rectMinusCorner_tileable_iff}
  \uses{LTileable, LTileable_swap_rectangleMinusCorner, rectMinusCorner_tileable_area_div_3, rectMinusCorner_tileable_of_area_m_le_n,
        rectangle, rectangleMinus2Corner, rectangleMinusCorner}
  \leanok
  \proven
  The L-tromino shape, anchored at origin -/ def lTromino : Prototile := ⟨[(0, 0), (0, 1), (1, 0)], by decide⟩ /-- A protoset containing just the L-tromino -/ def lTrominoSet : Protoset Unit := ⟨fun \_ => lTromino⟩ /- \#\# Utility Lemmas -/ /-- Each L-tromino covers exactly 3 cells -/ theorem lTromino...
\end{theorem}

\begin{definition}
  \label{lTrominoCellsAt}
  \lean{lTrominoCellsAt}
  \uses{Cell, PlacedTile, PlacedTile_cells, lTrominoSet}
  \leanok
  \proven
  The L-tromino shape, anchored at origin -/ def lTromino : Prototile := ⟨[(0, 0), (0, 1), (1, 0)], by decide⟩ /-- A protoset containing just the L-tromino -/ def lTrominoSet : Protoset Unit := ⟨fun \_ => lTromino⟩ /- \#\# Utility Lemmas -/ /-- Each L-tromino covers exactly 3 cells -/ theorem lTromino...
\end{definition}

\begin{definition}
  \label{singleLTrominoTiling}
  \lean{singleLTrominoTiling}
  \uses{TileSet, lTrominoSet}
  \leanok
  \proven
  The L-tromino shape, anchored at origin -/ def lTromino : Prototile := ⟨[(0, 0), (0, 1), (1, 0)], by decide⟩ /-- A protoset containing just the L-tromino -/ def lTrominoSet : Protoset Unit := ⟨fun \_ => lTromino⟩ /- \#\# Utility Lemmas -/ /-- Each L-tromino covers exactly 3 cells -/ theorem lTromino...
\end{definition}

\begin{theorem}
  \label{singleLTromino_valid}
  \lean{singleLTromino_valid}
  \uses{TileSet, Valid, cellsAt, coveredCells,
        lTrominoCellsAt, singleLTrominoTiling}
  \leanok
  \proven
  The L-tromino shape, anchored at origin -/ def lTromino : Prototile := ⟨[(0, 0), (0, 1), (1, 0)], by decide⟩ /-- A protoset containing just the L-tromino -/ def lTrominoSet : Protoset Unit := ⟨fun \_ => lTromino⟩ /- \#\# Utility Lemmas -/ /-- Each L-tromino covers exactly 3 cells -/ theorem lTromino...
\end{theorem}

\begin{theorem}
  \label{LTileable_single_lTromino}
  \lean{LTileable_single_lTromino}
  \uses{LTileable, lTrominoCellsAt, singleLTrominoTiling, singleLTromino_valid}
  \leanok
  \proven
  The L-tromino shape, anchored at origin -/ def lTromino : Prototile := ⟨[(0, 0), (0, 1), (1, 0)], by decide⟩ /-- A protoset containing just the L-tromino -/ def lTrominoSet : Protoset Unit := ⟨fun \_ => lTromino⟩ /- \#\# Utility Lemmas -/ /-- Each L-tromino covers exactly 3 cells -/ theorem lTromino...
\end{theorem}

\begin{definition}
  \label{topRightLTromino}
  \lean{topRightLTromino}
  \uses{Cell, lTrominoCellsAt}
  \leanok
  \proven
  The L-tromino shape, anchored at origin -/ def lTromino : Prototile := ⟨[(0, 0), (0, 1), (1, 0)], by decide⟩ /-- A protoset containing just the L-tromino -/ def lTrominoSet : Protoset Unit := ⟨fun \_ => lTromino⟩ /- \#\# Utility Lemmas -/ /-- Each L-tromino covers exactly 3 cells -/ theorem lTromino...
\end{definition}

\begin{theorem}
  \label{LTileable_topRightLTromino}
  \lean{LTileable_topRightLTromino}
  \uses{LTileable, LTileable_single_lTromino, rectangle, rectangleMinus2Corner,
        topRightLTromino}
  \leanok
  \proven
  The L-tromino shape, anchored at origin -/ def lTromino : Prototile := ⟨[(0, 0), (0, 1), (1, 0)], by decide⟩ /-- A protoset containing just the L-tromino -/ def lTrominoSet : Protoset Unit := ⟨fun \_ => lTromino⟩ /- \#\# Utility Lemmas -/ /-- Each L-tromino covers exactly 3 cells -/ theorem lTromino...
\end{theorem}

\section{RExp Representations for Decomposition}

\begin{definition}
  \label{rectangleMinus2Corner_rexp}
  \lean{rectangleMinus2Corner_rexp}
  \uses{RExp_r, topRightLTromino}
  \leanok
  \proven
  The L-tromino shape, anchored at origin -/ def lTromino : Prototile := ⟨[(0, 0), (0, 1), (1, 0)], by decide⟩ /-- A protoset containing just the L-tromino -/ def lTrominoSet : Protoset Unit := ⟨fun \_ => lTromino⟩ /- \#\# Utility Lemmas -/ /-- Each L-tromino covers exactly 3 cells -/ theorem lTromino...
\end{definition}

\begin{definition}
  \label{topRightLTromino_rexp}
  \lean{topRightLTromino_rexp}
  \uses{RExp_r, rectangle, translateRegion}
  \leanok
  \proven
  The L-tromino shape, anchored at origin -/ def lTromino : Prototile := ⟨[(0, 0), (0, 1), (1, 0)], by decide⟩ /-- A protoset containing just the L-tromino -/ def lTrominoSet : Protoset Unit := ⟨fun \_ => lTromino⟩ /- \#\# Utility Lemmas -/ /-- Each L-tromino covers exactly 3 cells -/ theorem lTromino...
\end{definition}

\begin{definition}
  \label{translateRectangle_rexp}
  \lean{translateRectangle_rexp}
  \uses{RExp_r}
  \leanok
  \proven
  The L-tromino shape, anchored at origin -/ def lTromino : Prototile := ⟨[(0, 0), (0, 1), (1, 0)], by decide⟩ /-- A protoset containing just the L-tromino -/ def lTrominoSet : Protoset Unit := ⟨fun \_ => lTromino⟩ /- \#\# Utility Lemmas -/ /-- Each L-tromino covers exactly 3 cells -/ theorem lTromino...
\end{definition}

\begin{theorem}
  \label{lTromino_rotation2_cells}
  \lean{lTromino_rotation2_cells}
  \uses{lTromino, rotateProto}
  \leanok
  \proven
  The L-tromino shape, anchored at origin -/ def lTromino : Prototile := ⟨[(0, 0), (0, 1), (1, 0)], by decide⟩ /-- A protoset containing just the L-tromino -/ def lTrominoSet : Protoset Unit := ⟨fun \_ => lTromino⟩ /- \#\# Utility Lemmas -/ /-- Each L-tromino covers exactly 3 cells -/ theorem lTromino...
\end{theorem}

\begin{theorem}
  \label{topRightLTromino_rexp_eval}
  \lean{topRightLTromino_rexp_eval}
  \uses{Cell, PlacedTile, PlacedTile_cells, RExp_eval,
        lTrominoCellsAt, lTrominoSet, lTromino_rotation2_cells, topRightLTromino,
        topRightLTromino_rexp, translateCell}
  \leanok
  \proven
  \((topRightLTromino\_rexp\_eval)\)
\end{theorem}

\begin{theorem}
  \label{rectangleMinus2Corner_rexp_eval}
  \lean{rectangleMinus2Corner_rexp_eval}
  \uses{Cell, RExp_eval, cornerTR, cornerTR2,
        mem_rectangle, rectangle, rectangleMinus2Corner, rectangleMinus2Corner_rexp,
        rectangleMinusCorner, topRightLTromino, translateRegion}
  \leanok
  \proven
  \((rectangleMinus2Corner\_rexp\_eval)\)
\end{theorem}

\begin{theorem}
  \label{rectangleMinus2Corner_decomp_rexp}
  \lean{rectangleMinus2Corner_decomp_rexp}
  \uses{RExp_eval, rectangle, rectangleMinus2Corner, rectangleMinus2Corner_rexp,
        rectangleMinusCorner, rectangleMinusCorner_rexp, topRightLTromino, topRightLTromino_rexp,
        translateRectangle_rexp, translateRegion}
  \leanok
  \proven
  The L-tromino shape, anchored at origin -/ def lTromino : Prototile := ⟨[(0, 0), (0, 1), (1, 0)], by decide⟩ /-- A protoset containing just the L-tromino -/ def lTrominoSet : Protoset Unit := ⟨fun \_ => lTromino⟩ /- \#\# Utility Lemmas -/ /-- Each L-tromino covers exactly 3 cells -/ theorem lTromino...
\end{theorem}

\begin{theorem}
  \label{rectangleMinus2Corner_decomp}
  \lean{rectangleMinus2Corner_decomp}
  \uses{Cell, RExp_eval, rect, rectangle,
        rectangleMinus2Corner, rectangleMinus2Corner_decomp_rexp, rectangleMinus2Corner_rexp, rectangleMinus2Corner_rexp_eval,
        rectangleMinusCorner, rectangleMinusCorner_rexp, rectangleMinusCorner_rexp_eval, rectangle_eq_rect_coe,
        topRightLTromino, topRightLTromino_rexp, topRightLTromino_rexp_eval, translateRectangle_rexp,
        translateRegion, translateRegion_eq_translate_coe}
  \leanok
  \proven
  The L-tromino shape, anchored at origin -/ def lTromino : Prototile := ⟨[(0, 0), (0, 1), (1, 0)], by decide⟩ /-- A protoset containing just the L-tromino -/ def lTrominoSet : Protoset Unit := ⟨fun \_ => lTromino⟩ /- \#\# Utility Lemmas -/ /-- Each L-tromino covers exactly 3 cells -/ theorem lTromino...
\end{theorem}

\begin{theorem}
  \label{rectangleMinus2Corner_decomp_disjoint_rexp_1}
  \lean{rectangleMinus2Corner_decomp_disjoint_rexp_1}
  \uses{RExp_eval, rectangleMinusCorner, rectangleMinusCorner_rexp, topRightLTromino,
        translateRectangle_rexp}
  \leanok
  \proven
  The L-tromino shape, anchored at origin -/ def lTromino : Prototile := ⟨[(0, 0), (0, 1), (1, 0)], by decide⟩ /-- A protoset containing just the L-tromino -/ def lTrominoSet : Protoset Unit := ⟨fun \_ => lTromino⟩ /- \#\# Utility Lemmas -/ /-- Each L-tromino covers exactly 3 cells -/ theorem lTromino...
\end{theorem}

\begin{theorem}
  \label{rectangleMinus2Corner_decomp_disjoint_rexp_2}
  \lean{rectangleMinus2Corner_decomp_disjoint_rexp_2}
  \uses{RExp_eval, rect, rectangleMinusCorner, rectangleMinusCorner_rexp,
        topRightLTromino, topRightLTromino_rexp}
  \leanok
  \proven
  The L-tromino shape, anchored at origin -/ def lTromino : Prototile := ⟨[(0, 0), (0, 1), (1, 0)], by decide⟩ /-- A protoset containing just the L-tromino -/ def lTrominoSet : Protoset Unit := ⟨fun \_ => lTromino⟩ /- \#\# Utility Lemmas -/ /-- Each L-tromino covers exactly 3 cells -/ theorem lTromino...
\end{theorem}

\begin{theorem}
  \label{rectangleMinus2Corner_decomp_disjoint_rexp_3}
  \lean{rectangleMinus2Corner_decomp_disjoint_rexp_3}
  \uses{RExp_eval, topRightLTromino, topRightLTromino_rexp, translateRectangle_rexp}
  \leanok
  \proven
  The L-tromino shape, anchored at origin -/ def lTromino : Prototile := ⟨[(0, 0), (0, 1), (1, 0)], by decide⟩ /-- A protoset containing just the L-tromino -/ def lTrominoSet : Protoset Unit := ⟨fun \_ => lTromino⟩ /- \#\# Utility Lemmas -/ /-- Each L-tromino covers exactly 3 cells -/ theorem lTromino...
\end{theorem}

\begin{theorem}
  \label{rectangleMinus2Corner_decomp_disjoint_rexp}
  \lean{rectangleMinus2Corner_decomp_disjoint_rexp}
  \uses{RExp_eval, rectangleMinus2Corner_decomp_disjoint_rexp_1, rectangleMinus2Corner_decomp_disjoint_rexp_2, rectangleMinus2Corner_decomp_disjoint_rexp_3,
        rectangleMinusCorner_rexp, topRightLTromino_rexp, translateRectangle_rexp}
  \leanok
  \proven
  The L-tromino shape, anchored at origin -/ def lTromino : Prototile := ⟨[(0, 0), (0, 1), (1, 0)], by decide⟩ /-- A protoset containing just the L-tromino -/ def lTrominoSet : Protoset Unit := ⟨fun \_ => lTromino⟩ /- \#\# Utility Lemmas -/ /-- Each L-tromino covers exactly 3 cells -/ theorem lTromino...
\end{theorem}

\begin{theorem}
  \label{rectangleMinus2Corner_decomp_disjoint}
  \lean{rectangleMinus2Corner_decomp_disjoint}
  \uses{Cell, RExp_eval, rect, rectangle,
        rectangleMinus2Corner_decomp_disjoint_rexp, rectangleMinusCorner, rectangleMinusCorner_rexp, rectangleMinusCorner_rexp_eval,
        rectangle_eq_rect_coe, topRightLTromino, topRightLTromino_rexp, topRightLTromino_rexp_eval,
        translateRectangle_rexp, translateRegion, translateRegion_eq_translate_coe}
  \leanok
  \proven
  The L-tromino shape, anchored at origin -/ def lTromino : Prototile := ⟨[(0, 0), (0, 1), (1, 0)], by decide⟩ /-- A protoset containing just the L-tromino -/ def lTrominoSet : Protoset Unit := ⟨fun \_ => lTromino⟩ /- \#\# Utility Lemmas -/ /-- Each L-tromino covers exactly 3 cells -/ theorem lTromino...
\end{theorem}

\begin{theorem}
  \label{tileable_rectangleMinus2Corner_3jplus2_3kplus1}
  \lean{tileable_rectangleMinus2Corner_3jplus2_3kplus1}
  \uses{LTileable, LTileable_single_lTromino, LTileable_translate, Tileable_union,
        lTrominoSet, rectMinusCorner_tileable_iff, rectangle, rectangleMinus2Corner,
        rectangleMinus2Corner_decomp, rectangleMinus2Corner_decomp_disjoint, rectangleMinusCorner, tileable_mult3_x2,
        topRightLTromino, translateRegion}
  \leanok
  \proven
  \((tileable\_rectangleMinus2Corner\_3jplus2\_3kplus1)\)
\end{theorem}

\begin{definition}
  \label{rotateRegion90}
  \lean{rotateRegion90}
  \uses{Region, rotateCell90}
  \leanok
  \proven
  \((rotateRegion90)\)
\end{definition}

\begin{theorem}
  \label{rotateCell90_rotateCell}
  \lean{rotateCell90_rotateCell}
  \uses{Cell, rotateCell, rotateCell90}
  \leanok
  \proven
  \((rotateCell90\_rotateCell)\)
\end{theorem}

\begin{theorem}
  \label{rotateProto_rotateCell90}
  \lean{rotateProto_rotateCell90}
  \uses{lTromino, rotateCell, rotateCell90, rotateCell90_rotateCell,
        rotateProto}
  \leanok
  \proven
  \((rotateProto\_rotateCell90)\)
\end{theorem}

\begin{theorem}
  \label{rotateCell90_translateCell}
  \lean{rotateCell90_translateCell}
  \uses{Cell, rotateCell90, translateCell}
  \leanok
  \proven
  \((rotateCell90\_translateCell)\)
\end{theorem}

\begin{theorem}
  \label{LTileable_rotate90}
  \lean{LTileable_rotate90}
  \uses{LTileable, PlacedTile, PlacedTile_cells, Region,
        TileSet, cellsAt, coveredCells, lTromino,
        lTrominoSet, rotateCell90, rotateCell90_translateCell, rotateProto,
        rotateProto_rotateCell90, rotateRegion90, translateCell}
  \leanok
  \proven
  \((LTileable\_rotate90)\)
\end{theorem}

\begin{definition}
  \label{piece1_rexp}
  \lean{piece1_rexp}
  \uses{RExp_r}
  \leanok
  \proven
  \((piece1\_rexp)\)
\end{definition}

\begin{definition}
  \label{piece2_rexp}
  \lean{piece2_rexp}
  \uses{RExp_r}
  \leanok
  \proven
  \((piece2\_rexp)\)
\end{definition}

\begin{definition}
  \label{piece1}
  \lean{piece1}
  \uses{Cell, lTrominoCellsAt}
  \leanok
  \proven
  \((piece1)\)
\end{definition}

\begin{definition}
  \label{piece2}
  \lean{piece2}
  \uses{Cell, rectangle}
  \leanok
  \proven
  \((piece2)\)
\end{definition}

\begin{theorem}
  \label{piece1_val_rexp}
  \lean{piece1_val_rexp}
  \uses{PlacedTile, PlacedTile_cells, lTromino, lTrominoCellsAt,
        lTrominoSet, piece1, piece1_rexp, rotateCell,
        rotateCell90, rotateProto, translateCell}
  \leanok
  \proven
  \((piece1\_val\_rexp)\)
\end{theorem}

\begin{theorem}
  \label{piece2_val_rexp}
  \lean{piece2_val_rexp}
  \uses{piece2, piece2_rexp, rectangle}
  \leanok
  \proven
  \((piece2\_val\_rexp)\)
\end{theorem}

\begin{theorem}
  \label{piece2_eq_rotate_rectangleMinusCorner}
  \lean{piece2_eq_rotate_rectangleMinusCorner}
  \uses{cornerTR, mem_rectangle, piece2, rectangleMinusCorner,
        rotateCell90, rotateRegion90, translateCell, translateRegion}
  \leanok
  \proven
  \((piece2\_eq\_rotate\_rectangleMinusCorner)\)
\end{theorem}

\begin{theorem}
  \label{tileable_piece2}
  \lean{tileable_piece2}
  \uses{LTileable, LTileable_rotate90, LTileable_translate, piece2,
        piece2_eq_rotate_rectangleMinusCorner, rectMinusCorner_tileable_iff}
  \leanok
  \proven
  \((tileable\_piece2)\)
\end{theorem}

\begin{theorem}
  \label{tileable_rectangleMinus2Corner_4_3kplus2}
  \lean{tileable_rectangleMinus2Corner_4_3kplus2}
  \uses{Cell, LTileable, LTileable_single_lTromino, Tileable_union,
        lTrominoSet, piece1, piece1_rexp, piece1_val_rexp,
        piece2, piece2_rexp, piece2_val_rexp, rectangleMinus2Corner,
        rectangleMinus2Corner_rexp, rectangleMinus2Corner_rexp_eval, tileable_piece2}
  \leanok
  \proven
  \((tileable\_rectangleMinus2Corner\_4\_3kplus2)\)
\end{theorem}

\begin{definition}
  \label{left_piece_j}
  \lean{left_piece_j}
  \uses{Cell, rectangle}
  \leanok
  \proven
  \((left\_piece\_j)\)
\end{definition}

\begin{definition}
  \label{right_piece_j}
  \lean{right_piece_j}
  \uses{Cell, rectangleMinus2Corner, translateRegion}
  \leanok
  \proven
  \((right\_piece\_j)\)
\end{definition}

\begin{theorem}
  \label{decomp_j}
  \lean{decomp_j}
  \uses{Cell, RExp_r, left_piece_j, rectangle,
        rectangleMinus2Corner, rectangleMinus2Corner_rexp, rectangleMinus2Corner_rexp_eval, right_piece_j,
        translate, translateCell, translateRegion}
  \leanok
  \proven
  \((decomp\_j)\)
\end{theorem}

\begin{theorem}
  \label{decomp_j_disj}
  \lean{decomp_j_disj}
  \uses{left_piece_j, rectangle, rectangleMinus2Corner_rexp, rectangleMinus2Corner_rexp_eval,
        right_piece_j, translateCell, translateRegion}
  \leanok
  \proven
  \((decomp\_j\_disj)\)
\end{theorem}

\begin{definition}
  \label{pieceA}
  \lean{pieceA}
  \uses{Cell, rectangleMinus2Corner, swapRegion}
  \leanok
  \proven
  \((pieceA)\)
\end{definition}

\begin{definition}
  \label{pieceA_rexp}
  \lean{pieceA_rexp}
  \uses{RExp_r}
  \leanok
  \proven
  \((pieceA\_rexp)\)
\end{definition}

\begin{theorem}
  \label{pieceA_val_rexp}
  \lean{pieceA_val_rexp}
  \uses{Cell, cornerTR, cornerTR2, mem_rectangle,
        pieceA, pieceA_rexp, rectangleMinus2Corner, swapRegion}
  \leanok
  \proven
  \((pieceA\_val\_rexp)\)
\end{theorem}

\begin{definition}
  \label{pieceB}
  \lean{pieceB}
  \uses{Cell, lTrominoCellsAt}
  \leanok
  \proven
  \((pieceB)\)
\end{definition}

\begin{definition}
  \label{pieceB_rexp}
  \lean{pieceB_rexp}
  \uses{RExp_r}
  \leanok
  \proven
  \((pieceB\_rexp)\)
\end{definition}

\begin{theorem}
  \label{pieceB_val_rexp}
  \lean{pieceB_val_rexp}
  \uses{Cell, PlacedTile, PlacedTile_cells, lTromino,
        lTrominoCellsAt, lTrominoSet, pieceB, pieceB_rexp,
        rotateCell, rotateCell90, rotateProto, translateCell}
  \leanok
  \proven
  \((pieceB\_val\_rexp)\)
\end{theorem}

\begin{definition}
  \label{pieceC}
  \lean{pieceC}
  \uses{Cell, rectangle, translateRegion}
  \leanok
  \proven
  \((pieceC)\)
\end{definition}

\begin{definition}
  \label{pieceC_rexp}
  \lean{pieceC_rexp}
  \uses{RExp_r}
  \leanok
  \proven
  \((pieceC\_rexp)\)
\end{definition}

\begin{theorem}
  \label{pieceC_val_rexp}
  \lean{pieceC_val_rexp}
  \uses{Cell, pieceC, pieceC_rexp, rect,
        rectangle, translateCell, translateRegion}
  \leanok
  \proven
  \((pieceC\_val\_rexp)\)
\end{theorem}

\begin{theorem}
  \label{decomp_7_rexp}
  \lean{decomp_7_rexp}
  \uses{pieceA_rexp, pieceB_rexp, pieceC_rexp, rect,
        rectangleMinus2Corner_rexp}
  \leanok
  \proven
  \((decomp\_7\_rexp)\)
\end{theorem}

\begin{theorem}
  \label{decomp_7}
  \lean{decomp_7}
  \uses{decomp_7_rexp, pieceA, pieceA_val_rexp, pieceB,
        pieceB_val_rexp, pieceC, pieceC_val_rexp, rectangleMinus2Corner,
        rectangleMinus2Corner_rexp_eval}
  \leanok
  \proven
  \((decomp\_7)\)
\end{theorem}

\begin{theorem}
  \label{decomp_7_disj1}
  \lean{decomp_7_disj1}
  \uses{pieceA, pieceA_rexp, pieceA_val_rexp, pieceB,
        pieceB_rexp, pieceB_val_rexp, rect}
  \leanok
  \proven
  \((decomp\_7\_disj1)\)
\end{theorem}

\begin{theorem}
  \label{decomp_7_disj2}
  \lean{decomp_7_disj2}
  \uses{pieceA, pieceA_rexp, pieceA_val_rexp, pieceB,
        pieceB_rexp, pieceB_val_rexp, pieceC, pieceC_rexp,
        pieceC_val_rexp, rect}
  \leanok
  \proven
  \((decomp\_7\_disj2)\)
\end{theorem}

\begin{theorem}
  \label{tileable_pieceA}
  \lean{tileable_pieceA}
  \uses{LTileable, LTileable_swap, pieceA, rectangleMinus2Corner,
        tileable_rectangleMinus2Corner_3jplus2_3kplus1}
  \leanok
  \proven
  \((tileable\_pieceA)\)
\end{theorem}

\begin{theorem}
  \label{tileable_pieceB}
  \lean{tileable_pieceB}
  \uses{LTileable, LTileable_single_lTromino, pieceB}
  \leanok
  \proven
  \((tileable\_pieceB)\)
\end{theorem}

\begin{theorem}
  \label{tileable_pieceC}
  \lean{tileable_pieceC}
  \uses{LTileable, LTileable_translate, pieceC, rect_tileable_iff}
  \leanok
  \proven
  \((tileable\_pieceC)\)
\end{theorem}

\begin{theorem}
  \label{tileable_rectangleMinus2Corner_3jplus1_3kplus2}
  \lean{tileable_rectangleMinus2Corner_3jplus1_3kplus2}
  \uses{LTileable, LTileable_translate, Tileable_union, decomp_7,
        decomp_7_disj1, decomp_7_disj2, decomp_j, decomp_j_disj,
        lTrominoSet, left_piece_j, pieceA, pieceB,
        pieceC, rect_tileable_iff, rectangle, rectangleMinus2Corner,
        right_piece_j, tileable_pieceA, tileable_pieceB, tileable_pieceC,
        tileable_rectangleMinus2Corner_4_3kplus2, translateRegion}
  \leanok
  \proven
  \((tileable\_rectangleMinus2Corner\_3jplus1\_3kplus2)\)
\end{theorem}

\begin{theorem}
  \label{rectMinus2Corner_tileable_of_area_mod2}
  \lean{rectMinus2Corner_tileable_of_area_mod2}
  \uses{LTileable, mod3_one_two_or_two_one_of_area_mod3_two, rectangleMinus2Corner, tileable_rectangleMinus2Corner_3jplus1_3kplus2,
        tileable_rectangleMinus2Corner_3jplus2_3kplus1}
  \leanok
  \proven
  The L-tromino shape, anchored at origin -/ def lTromino : Prototile := ⟨[(0, 0), (0, 1), (1, 0)], by decide⟩ /-- A protoset containing just the L-tromino -/ def lTrominoSet : Protoset Unit := ⟨fun \_ => lTromino⟩ /- \#\# Utility Lemmas -/ /-- Each L-tromino covers exactly 3 cells -/ theorem lTromino...
\end{theorem}

\chapter{Rectangle and Omega Infrastructure}

\begin{theorem}
  \label{example}
  \lean{example}
  \uses{Cell, rect, rectangle}
  \leanok
  \proven
  \((example)\)
\end{theorem}

\section{Basic Definitions}

\begin{definition}
  \label{rect}
  \lean{rect}
  \uses{Cell}
  \leanok
  \proven
  Half-open rectangle: `rect x0 y0 x1 y1` contains all cells `(x, y)` such that `x0 ≤ x < x1` and `y0 ≤ y < y1`.
\end{definition}

\begin{definition}
  \label{translate}
  \lean{translate}
  \uses{Cell, TileSet}
  \leanok
  \proven
  A cell on the integer grid -/ abbrev Cell := ℤ × ℤ /-- A region is a finite set of cells -/ abbrev Region := Finset Cell /- \#\# Prototile and Protoset -/ /- A prototile is the "shape" of a tile, represented as a list of cells with no duplicates. This design allows both: * Computation: iterate over...
\end{definition}

\begin{definition}
  \label{rotate}
  \lean{rotate}
  \uses{Cell, rotateCell}
  \leanok
  \proven
  Half-open rectangle: `rect x0 y0 x1 y1` contains all cells `(x, y)` such that `x0 ≤ x < x1` and `y0 ≤ y < y1`. -/ def rect (x0 y0 x1 y1 : ℤ) : Set Cell := \{p : Cell | x0 ≤ p.1 ∧ p.1 < x1 ∧ y0 ≤ p.2 ∧ p.2 < y1\} /-- Translation of a set of cells by `(dx, dy)` -/ def translate (dx dy : ℤ) (s : Set C...
\end{definition}

\begin{definition}
  \label{reflect}
  \lean{reflect}
  \uses{Cell, rect, rotate, rotateCell,
        swapCell, translate}
  \leanok
  \proven
  Half-open rectangle: `rect x0 y0 x1 y1` contains all cells `(x, y)` such that `x0 ≤ x < x1` and `y0 ≤ y < y1`. -/ def rect (x0 y0 x1 y1 : ℤ) : Set Cell := \{p : Cell | x0 ≤ p.1 ∧ p.1 < x1 ∧ y0 ≤ p.2 ∧ p.2 < y1\} /-- Translation of a set of cells by `(dx, dy)` -/ def translate (dx dy : ℤ) (s : Set C...
\end{definition}

\section{Simplification Lemmas}

\begin{theorem}
  \label{swapCell_def}
  \lean{swapCell_def}
  \uses{Cell, rectangle, rotateCell, rotateCell90,
        swapCell}
  \leanok
  \proven
  Half-open rectangle: `rect x0 y0 x1 y1` contains all cells `(x, y)` such that `x0 ≤ x < x1` and `y0 ≤ y < y1`. -/ def rect (x0 y0 x1 y1 : ℤ) : Set Cell := \{p : Cell | x0 ≤ p.1 ∧ p.1 < x1 ∧ y0 ≤ p.2 ∧ p.2 < y1\} /-- Translation of a set of cells by `(dx, dy)` -/ def translate (dx dy : ℤ) (s : Set C...
\end{theorem}

\section{Finite Boxes}

\begin{definition}
  \label{boxFinset}
  \lean{boxFinset}
  \uses{Cell, rect}
  \leanok
  \proven
  Half-open rectangle: `rect x0 y0 x1 y1` contains all cells `(x, y)` such that `x0 ≤ x < x1` and `y0 ≤ y < y1`. -/ def rect (x0 y0 x1 y1 : ℤ) : Set Cell := \{p : Cell | x0 ≤ p.1 ∧ p.1 < x1 ∧ y0 ≤ p.2 ∧ p.2 < y1\} /-- Translation of a set of cells by `(dx, dy)` -/ def translate (dx dy : ℤ) (s : Set C...
\end{definition}

\begin{theorem}
  \label{card_boxFinset}
  \lean{card_boxFinset}
  \uses{boxFinset, example}
  \leanok
  \proven
  Half-open rectangle: `rect x0 y0 x1 y1` contains all cells `(x, y)` such that `x0 ≤ x < x1` and `y0 ≤ y < y1`. -/ def rect (x0 y0 x1 y1 : ℤ) : Set Cell := \{p : Cell | x0 ≤ p.1 ∧ p.1 < x1 ∧ y0 ≤ p.2 ∧ p.2 < y1\} /-- Translation of a set of cells by `(dx, dy)` -/ def translate (dx dy : ℤ) (s : Set C...
\end{theorem}

\begin{theorem}
  \label{rect_eq_boxFinset_coe}
  \lean{rect_eq_boxFinset_coe}
  \uses{Cell, boxFinset, rect}
  \leanok
  \proven
  \((rect\_eq\_boxFinset\_coe)\)
\end{theorem}

\begin{definition}
  \label{rectFinset}
  \lean{rectFinset}
  \uses{Cell, boxFinset}
  \leanok
  \proven
  \((rectFinset)\)
\end{definition}

\begin{theorem}
  \label{card_rectFinset}
  \lean{card_rectFinset}
  \uses{card_boxFinset, rectFinset}
  \leanok
  \proven
  \((card\_rectFinset)\)
\end{theorem}

\section{Cardinality Lemmas}

\begin{theorem}
  \label{rect_encard}
  \lean{rect_encard}
  \uses{boxFinset, rect, rect_eq_boxFinset_coe}
  \leanok
  \proven
  Half-open rectangle: `rect x0 y0 x1 y1` contains all cells `(x, y)` such that `x0 ≤ x < x1` and `y0 ≤ y < y1`. -/ def rect (x0 y0 x1 y1 : ℤ) : Set Cell := \{p : Cell | x0 ≤ p.1 ∧ p.1 < x1 ∧ y0 ≤ p.2 ∧ p.2 < y1\} /-- Translation of a set of cells by `(dx, dy)` -/ def translate (dx dy : ℤ) (s : Set C...
\end{theorem}

\begin{theorem}
  \label{rect_encard_formula}
  \lean{rect_encard_formula}
  \uses{card_boxFinset, rect, rect_encard}
  \leanok
  \proven
  \((rect\_encard\_formula)\)
\end{theorem}

\begin{theorem}
  \label{rect_finite}
  \lean{rect_finite}
  \uses{rect, rect_eq_boxFinset_coe, rectangle}
  \leanok
  \proven
  Half-open rectangle: `rect x0 y0 x1 y1` contains all cells `(x, y)` such that `x0 ≤ x < x1` and `y0 ≤ y < y1`. -/ def rect (x0 y0 x1 y1 : ℤ) : Set Cell := \{p : Cell | x0 ≤ p.1 ∧ p.1 < x1 ∧ y0 ≤ p.2 ∧ p.2 < y1\} /-- Translation of a set of cells by `(dx, dy)` -/ def translate (dx dy : ℤ) (s : Set C...
\end{theorem}

\begin{theorem}
  \label{rect_ncard}
  \lean{rect_ncard}
  \uses{Cell, card_boxFinset, rect, rect_eq_boxFinset_coe}
  \leanok
  \proven
  Half-open rectangle: `rect x0 y0 x1 y1` contains all cells `(x, y)` such that `x0 ≤ x < x1` and `y0 ≤ y < y1`. -/ def rect (x0 y0 x1 y1 : ℤ) : Set Cell := \{p : Cell | x0 ≤ p.1 ∧ p.1 < x1 ∧ y0 ≤ p.2 ∧ p.2 < y1\} /-- Translation of a set of cells by `(dx, dy)` -/ def translate (dx dy : ℤ) (s : Set C...
\end{theorem}

\begin{theorem}
  \label{union_ncard}
  \lean{union_ncard}
  \uses{Cell}
  \leanok
  \proven
  Half-open rectangle: `rect x0 y0 x1 y1` contains all cells `(x, y)` such that `x0 ≤ x < x1` and `y0 ≤ y < y1`. -/ def rect (x0 y0 x1 y1 : ℤ) : Set Cell := \{p : Cell | x0 ≤ p.1 ∧ p.1 < x1 ∧ y0 ≤ p.2 ∧ p.2 < y1\} /-- Translation of a set of cells by `(dx, dy)` -/ def translate (dx dy : ℤ) (s : Set C...
\end{theorem}

\begin{theorem}
  \label{union_ncard_disjoint}
  \lean{union_ncard_disjoint}
  \uses{Cell, union_ncard}
  \leanok
  \proven
  Half-open rectangle: `rect x0 y0 x1 y1` contains all cells `(x, y)` such that `x0 ≤ x < x1` and `y0 ≤ y < y1`. -/ def rect (x0 y0 x1 y1 : ℤ) : Set Cell := \{p : Cell | x0 ≤ p.1 ∧ p.1 < x1 ∧ y0 ≤ p.2 ∧ p.2 < y1\} /-- Translation of a set of cells by `(dx, dy)` -/ def translate (dx dy : ℤ) (s : Set C...
\end{theorem}

\begin{theorem}
  \label{diff_ncard}
  \lean{diff_ncard}
  \uses{Cell, union_ncard_disjoint}
  \leanok
  \proven
  Half-open rectangle: `rect x0 y0 x1 y1` contains all cells `(x, y)` such that `x0 ≤ x < x1` and `y0 ≤ y < y1`. -/ def rect (x0 y0 x1 y1 : ℤ) : Set Cell := \{p : Cell | x0 ≤ p.1 ∧ p.1 < x1 ∧ y0 ≤ p.2 ∧ p.2 < y1\} /-- Translation of a set of cells by `(dx, dy)` -/ def translate (dx dy : ℤ) (s : Set C...
\end{theorem}

\begin{theorem}
  \label{diff_ncard_disjoint}
  \lean{diff_ncard_disjoint}
  \uses{Cell, diff_ncard}
  \leanok
  \proven
  Half-open rectangle: `rect x0 y0 x1 y1` contains all cells `(x, y)` such that `x0 ≤ x < x1` and `y0 ≤ y < y1`. -/ def rect (x0 y0 x1 y1 : ℤ) : Set Cell := \{p : Cell | x0 ≤ p.1 ∧ p.1 < x1 ∧ y0 ≤ p.2 ∧ p.2 < y1\} /-- Translation of a set of cells by `(dx, dy)` -/ def translate (dx dy : ℤ) (s : Set C...
\end{theorem}

\begin{theorem}
  \label{translate_ncard}
  \lean{translate_ncard}
  \uses{Cell, translate, translateCell, translateCell_injective}
  \leanok
  \proven
  Half-open rectangle: `rect x0 y0 x1 y1` contains all cells `(x, y)` such that `x0 ≤ x < x1` and `y0 ≤ y < y1`. -/ def rect (x0 y0 x1 y1 : ℤ) : Set Cell := \{p : Cell | x0 ≤ p.1 ∧ p.1 < x1 ∧ y0 ≤ p.2 ∧ p.2 < y1\} /-- Translation of a set of cells by `(dx, dy)` -/ def translate (dx dy : ℤ) (s : Set C...
\end{theorem}

\begin{theorem}
  \label{rotate_ncard}
  \lean{rotate_ncard}
  \uses{Cell, rotate, rotateCell, rotateCell_injective}
  \leanok
  \proven
  Half-open rectangle: `rect x0 y0 x1 y1` contains all cells `(x, y)` such that `x0 ≤ x < x1` and `y0 ≤ y < y1`. -/ def rect (x0 y0 x1 y1 : ℤ) : Set Cell := \{p : Cell | x0 ≤ p.1 ∧ p.1 < x1 ∧ y0 ≤ p.2 ∧ p.2 < y1\} /-- Translation of a set of cells by `(dx, dy)` -/ def translate (dx dy : ℤ) (s : Set C...
\end{theorem}

\begin{theorem}
  \label{reflect_ncard}
  \lean{reflect_ncard}
  \uses{Cell, rect, rectangle, reflect,
        rotate, swapCell, swapCell_def, swapCell_injective,
        swapCell_involutive}
  \leanok
  \proven
  Half-open rectangle: `rect x0 y0 x1 y1` contains all cells `(x, y)` such that `x0 ≤ x < x1` and `y0 ≤ y < y1`. -/ def rect (x0 y0 x1 y1 : ℤ) : Set Cell := \{p : Cell | x0 ≤ p.1 ∧ p.1 < x1 ∧ y0 ≤ p.2 ∧ p.2 < y1\} /-- Translation of a set of cells by `(dx, dy)` -/ def translate (dx dy : ℤ) (s : Set C...
\end{theorem}

\section{Automation Tactic}

\section{Rectangle Expression DSL}

\begin{definition}
  \label{RExp_eval}
  \lean{RExp.eval}
  \uses{Cell, rect, reflect, rotate,
        translate}
  \leanok
  \proven
  Half-open rectangle: `rect x0 y0 x1 y1` contains all cells `(x, y)` such that `x0 ≤ x < x1` and `y0 ≤ y < y1`. -/ def rect (x0 y0 x1 y1 : ℤ) : Set Cell := \{p : Cell | x0 ≤ p.1 ∧ p.1 < x1 ∧ y0 ≤ p.2 ∧ p.2 < y1\} /-- Translation of a set of cells by `(dx, dy)` -/ def translate (dx dy : ℤ) (s : Set C...
\end{definition}

\begin{definition}
  \label{RExp_r}
  \lean{RExp.r}
  \uses{Cell, RExp_eval, rect, rectangle,
        reflect, rotate, translate}
  \leanok
  \proven
  Half-open rectangle: `rect x0 y0 x1 y1` contains all cells `(x, y)` such that `x0 ≤ x < x1` and `y0 ≤ y < y1`. -/ def rect (x0 y0 x1 y1 : ℤ) : Set Cell := \{p : Cell | x0 ≤ p.1 ∧ p.1 < x1 ∧ y0 ≤ p.2 ∧ p.2 < y1\} /-- Translation of a set of cells by `(dx, dy)` -/ def translate (dx dy : ℤ) (s : Set C...
\end{definition}

\section{RExp Simplification Lemmas}

\section{RExp Cardinality Lemmas}

\begin{theorem}
  \label{RExp_ncard_rect}
  \lean{RExp.ncard_rect}
  \uses{RExp_eval, rect, rect_ncard}
  \leanok
  \proven
  Half-open rectangle: `rect x0 y0 x1 y1` contains all cells `(x, y)` such that `x0 ≤ x < x1` and `y0 ≤ y < y1`. -/ def rect (x0 y0 x1 y1 : ℤ) : Set Cell := \{p : Cell | x0 ≤ p.1 ∧ p.1 < x1 ∧ y0 ≤ p.2 ∧ p.2 < y1\} /-- Translation of a set of cells by `(dx, dy)` -/ def translate (dx dy : ℤ) (s : Set C...
\end{theorem}

\begin{theorem}
  \label{RExp_ncard_union}
  \lean{RExp.ncard_union}
  \uses{RExp_eval, union_ncard}
  \leanok
  \proven
  Half-open rectangle: `rect x0 y0 x1 y1` contains all cells `(x, y)` such that `x0 ≤ x < x1` and `y0 ≤ y < y1`. -/ def rect (x0 y0 x1 y1 : ℤ) : Set Cell := \{p : Cell | x0 ≤ p.1 ∧ p.1 < x1 ∧ y0 ≤ p.2 ∧ p.2 < y1\} /-- Translation of a set of cells by `(dx, dy)` -/ def translate (dx dy : ℤ) (s : Set C...
\end{theorem}

\begin{theorem}
  \label{RExp_ncard_union_disjoint}
  \lean{RExp.ncard_union_disjoint}
  \uses{RExp_eval, union_ncard_disjoint}
  \leanok
  \proven
  Half-open rectangle: `rect x0 y0 x1 y1` contains all cells `(x, y)` such that `x0 ≤ x < x1` and `y0 ≤ y < y1`. -/ def rect (x0 y0 x1 y1 : ℤ) : Set Cell := \{p : Cell | x0 ≤ p.1 ∧ p.1 < x1 ∧ y0 ≤ p.2 ∧ p.2 < y1\} /-- Translation of a set of cells by `(dx, dy)` -/ def translate (dx dy : ℤ) (s : Set C...
\end{theorem}

\begin{theorem}
  \label{RExp_ncard_inter}
  \lean{RExp.ncard_inter}
  \uses{RExp_eval}
  \leanok
  \proven
  Half-open rectangle: `rect x0 y0 x1 y1` contains all cells `(x, y)` such that `x0 ≤ x < x1` and `y0 ≤ y < y1`. -/ def rect (x0 y0 x1 y1 : ℤ) : Set Cell := \{p : Cell | x0 ≤ p.1 ∧ p.1 < x1 ∧ y0 ≤ p.2 ∧ p.2 < y1\} /-- Translation of a set of cells by `(dx, dy)` -/ def translate (dx dy : ℤ) (s : Set C...
\end{theorem}

\begin{theorem}
  \label{RExp_ncard_diff}
  \lean{RExp.ncard_diff}
  \uses{RExp_eval, diff_ncard}
  \leanok
  \proven
  Half-open rectangle: `rect x0 y0 x1 y1` contains all cells `(x, y)` such that `x0 ≤ x < x1` and `y0 ≤ y < y1`. -/ def rect (x0 y0 x1 y1 : ℤ) : Set Cell := \{p : Cell | x0 ≤ p.1 ∧ p.1 < x1 ∧ y0 ≤ p.2 ∧ p.2 < y1\} /-- Translation of a set of cells by `(dx, dy)` -/ def translate (dx dy : ℤ) (s : Set C...
\end{theorem}

\begin{theorem}
  \label{RExp_ncard_diff_disjoint}
  \lean{RExp.ncard_diff_disjoint}
  \uses{RExp_eval, diff_ncard_disjoint}
  \leanok
  \proven
  Half-open rectangle: `rect x0 y0 x1 y1` contains all cells `(x, y)` such that `x0 ≤ x < x1` and `y0 ≤ y < y1`. -/ def rect (x0 y0 x1 y1 : ℤ) : Set Cell := \{p : Cell | x0 ≤ p.1 ∧ p.1 < x1 ∧ y0 ≤ p.2 ∧ p.2 < y1\} /-- Translation of a set of cells by `(dx, dy)` -/ def translate (dx dy : ℤ) (s : Set C...
\end{theorem}

\begin{theorem}
  \label{RExp_ncard_shift}
  \lean{RExp.ncard_shift}
  \uses{RExp_eval, translate_ncard}
  \leanok
  \proven
  Half-open rectangle: `rect x0 y0 x1 y1` contains all cells `(x, y)` such that `x0 ≤ x < x1` and `y0 ≤ y < y1`. -/ def rect (x0 y0 x1 y1 : ℤ) : Set Cell := \{p : Cell | x0 ≤ p.1 ∧ p.1 < x1 ∧ y0 ≤ p.2 ∧ p.2 < y1\} /-- Translation of a set of cells by `(dx, dy)` -/ def translate (dx dy : ℤ) (s : Set C...
\end{theorem}

\begin{theorem}
  \label{RExp_ncard_rotate}
  \lean{RExp.ncard_rotate}
  \uses{RExp_eval, rotate, rotate_ncard}
  \leanok
  \proven
  Half-open rectangle: `rect x0 y0 x1 y1` contains all cells `(x, y)` such that `x0 ≤ x < x1` and `y0 ≤ y < y1`. -/ def rect (x0 y0 x1 y1 : ℤ) : Set Cell := \{p : Cell | x0 ≤ p.1 ∧ p.1 < x1 ∧ y0 ≤ p.2 ∧ p.2 < y1\} /-- Translation of a set of cells by `(dx, dy)` -/ def translate (dx dy : ℤ) (s : Set C...
\end{theorem}

\begin{theorem}
  \label{RExp_ncard_reflect}
  \lean{RExp.ncard_reflect}
  \uses{Cell, RExp_eval, RExp_ncard_diff, RExp_ncard_diff_disjoint,
        RExp_ncard_inter, RExp_ncard_rect, RExp_ncard_rotate, RExp_ncard_shift,
        RExp_ncard_union, RExp_ncard_union_disjoint, RExp_r, diff_ncard,
        diff_ncard_disjoint, example, rect, rect_finite,
        rect_ncard, rectangle, reflect, reflect_ncard,
        swapCell_def, union_ncard, union_ncard_disjoint}
  \leanok
  \proven
  Half-open rectangle: `rect x0 y0 x1 y1` contains all cells `(x, y)` such that `x0 ≤ x < x1` and `y0 ≤ y < y1`. -/ def rect (x0 y0 x1 y1 : ℤ) : Set Cell := \{p : Cell | x0 ≤ p.1 ∧ p.1 < x1 ∧ y0 ≤ p.2 ∧ p.2 < y1\} /-- Translation of a set of cells by `(dx, dy)` -/ def translate (dx dy : ℤ) (s : Set C...
\end{theorem}

\section{RExp Automation Tactic}

\section{RExp Cardinality Automation Tactic}

\section{Demo Theorems}

\begin{theorem}
  \label{rect_left_right_disjoint}
  \lean{rect_left_right_disjoint}
  \uses{rect}
  \leanok
  \proven
  Half-open rectangle: `rect x0 y0 x1 y1` contains all cells `(x, y)` such that `x0 ≤ x < x1` and `y0 ≤ y < y1`. -/ def rect (x0 y0 x1 y1 : ℤ) : Set Cell := \{p : Cell | x0 ≤ p.1 ∧ p.1 < x1 ∧ y0 ≤ p.2 ∧ p.2 < y1\} /-- Translation of a set of cells by `(dx, dy)` -/ def translate (dx dy : ℤ) (s : Set C...
\end{theorem}

\begin{theorem}
  \label{translate_rect}
  \lean{translate_rect}
  \uses{RExp_eval, RExp_ncard_reflect, RExp_ncard_rotate, RExp_ncard_shift,
        RExp_r, example, rect, rect_finite,
        rect_ncard, rectangle, rectangleMinus2Corner_decomp_rexp, translate}
  \leanok
  \proven
  Half-open rectangle: `rect x0 y0 x1 y1` contains all cells `(x, y)` such that `x0 ≤ x < x1` and `y0 ≤ y < y1`. -/ def rect (x0 y0 x1 y1 : ℤ) : Set Cell := \{p : Cell | x0 ≤ p.1 ∧ p.1 < x1 ∧ y0 ≤ p.2 ∧ p.2 < y1\} /-- Translation of a set of cells by `(dx, dy)` -/ def translate (dx dy : ℤ) (s : Set C...
\end{theorem}

\section{RExp Example Theorems}

\section{Cardinality Example Theorems}

\chapter{Tiling Infrastructure}

\section{Basic Definitions}

\begin{definition}
  \label{Cell}
  \lean{Cell}
  \leanok
  \proven
  A cell on the integer grid
\end{definition}

\begin{definition}
  \label{Region}
  \lean{Region}
  \uses{Cell, Protoset, Prototile}
  \leanok
  \proven
  A cell on the integer grid -/ abbrev Cell := ℤ × ℤ /-- A region is a finite set of cells
\end{definition}

\section{Prototile and Protoset}

\begin{definition}
  \label{Prototile}
  \lean{Prototile}
  \uses{Cell}
  \leanok
  \proven
  A cell on the integer grid -/ abbrev Cell := ℤ × ℤ /-- A region is a finite set of cells -/ abbrev Region := Finset Cell /- \#\# Prototile and Protoset -/ /- A prototile is the "shape" of a tile, represented as a list of cells with no duplicates. This design allows both: * Computation: iterate over...
\end{definition}

\begin{theorem}
  \label{Prototile_coe_def}
  \lean{Prototile.coe_def}
  \uses{Cell, Prototile}
  \leanok
  \proven
  \((Prototile.coe\_def)\)
\end{theorem}

\begin{theorem}
  \label{Prototile_mem_coe}
  \lean{Prototile.mem_coe}
  \uses{Cell, Prototile, Prototile_coe_def}
  \leanok
  \proven
  \((Prototile.mem\_coe)\)
\end{theorem}

\begin{definition}
  \label{Prototile_card}
  \lean{Prototile.card}
  \uses{Prototile}
  \leanok
  \proven
  A cell on the integer grid -/ abbrev Cell := ℤ × ℤ /-- A region is a finite set of cells -/ abbrev Region := Finset Cell /- \#\# Prototile and Protoset -/ /- A prototile is the "shape" of a tile, represented as a list of cells with no duplicates. This design allows both: * Computation: iterate over...
\end{definition}

\begin{theorem}
  \label{Prototile_card_eq}
  \lean{Prototile.card_eq}
  \uses{Cell, Prototile, Prototile_card}
  \leanok
  \proven
  \((Prototile.card\_eq)\)
\end{theorem}

\begin{definition}
  \label{Protoset}
  \lean{Protoset}
  \uses{Prototile}
  \leanok
  \proven
  A cell on the integer grid -/ abbrev Cell := ℤ × ℤ /-- A region is a finite set of cells -/ abbrev Region := Finset Cell /- \#\# Prototile and Protoset -/ /- A prototile is the "shape" of a tile, represented as a list of cells with no duplicates. This design allows both: * Computation: iterate over...
\end{definition}

\section{Rotations}

\begin{definition}
  \label{rotateCell90}
  \lean{rotateCell90}
  \uses{Cell}
  \leanok
  \proven
  A cell on the integer grid -/ abbrev Cell := ℤ × ℤ /-- A region is a finite set of cells -/ abbrev Region := Finset Cell /- \#\# Prototile and Protoset -/ /- A prototile is the "shape" of a tile, represented as a list of cells with no duplicates. This design allows both: * Computation: iterate over...
\end{definition}

\begin{definition}
  \label{rotateCell}
  \lean{rotateCell}
  \uses{Cell, rotateCell90}
  \leanok
  \proven
  A cell on the integer grid -/ abbrev Cell := ℤ × ℤ /-- A region is a finite set of cells -/ abbrev Region := Finset Cell /- \#\# Prototile and Protoset -/ /- A prototile is the "shape" of a tile, represented as a list of cells with no duplicates. This design allows both: * Computation: iterate over...
\end{definition}

\begin{theorem}
  \label{rotateCell_injective}
  \lean{rotateCell_injective}
  \uses{rotateCell, rotateCell90}
  \leanok
  \proven
  A cell on the integer grid -/ abbrev Cell := ℤ × ℤ /-- A region is a finite set of cells -/ abbrev Region := Finset Cell /- \#\# Prototile and Protoset -/ /- A prototile is the "shape" of a tile, represented as a list of cells with no duplicates. This design allows both: * Computation: iterate over...
\end{theorem}

\begin{definition}
  \label{rotateProto}
  \lean{rotateProto}
  \uses{Cell, Prototile, rotateCell}
  \leanok
  \proven
  A cell on the integer grid -/ abbrev Cell := ℤ × ℤ /-- A region is a finite set of cells -/ abbrev Region := Finset Cell /- \#\# Prototile and Protoset -/ /- A prototile is the "shape" of a tile, represented as a list of cells with no duplicates. This design allows both: * Computation: iterate over...
\end{definition}

\begin{theorem}
  \label{rotateProto_card}
  \lean{rotateProto_card}
  \uses{Prototile, rotateCell_injective, rotateProto}
  \leanok
  \proven
  A cell on the integer grid -/ abbrev Cell := ℤ × ℤ /-- A region is a finite set of cells -/ abbrev Region := Finset Cell /- \#\# Prototile and Protoset -/ /- A prototile is the "shape" of a tile, represented as a list of cells with no duplicates. This design allows both: * Computation: iterate over...
\end{theorem}

\section{Placed Tiles}

\begin{definition}
  \label{PlacedTile}
  \lean{PlacedTile}
  \uses{Cell}
  \leanok
  \proven
  A cell on the integer grid -/ abbrev Cell := ℤ × ℤ /-- A region is a finite set of cells -/ abbrev Region := Finset Cell /- \#\# Prototile and Protoset -/ /- A prototile is the "shape" of a tile, represented as a list of cells with no duplicates. This design allows both: * Computation: iterate over...
\end{definition}

\begin{definition}
  \label{translateCell}
  \lean{translateCell}
  \uses{Cell}
  \leanok
  \proven
  A cell on the integer grid -/ abbrev Cell := ℤ × ℤ /-- A region is a finite set of cells -/ abbrev Region := Finset Cell /- \#\# Prototile and Protoset -/ /- A prototile is the "shape" of a tile, represented as a list of cells with no duplicates. This design allows both: * Computation: iterate over...
\end{definition}

\begin{theorem}
  \label{translateCell_injective}
  \lean{translateCell_injective}
  \uses{Cell, translateCell}
  \leanok
  \proven
  A cell on the integer grid -/ abbrev Cell := ℤ × ℤ /-- A region is a finite set of cells -/ abbrev Region := Finset Cell /- \#\# Prototile and Protoset -/ /- A prototile is the "shape" of a tile, represented as a list of cells with no duplicates. This design allows both: * Computation: iterate over...
\end{theorem}

\begin{definition}
  \label{translateRegion}
  \lean{translateRegion}
  \uses{Cell, Region, translateCell}
  \leanok
  \proven
  A cell on the integer grid -/ abbrev Cell := ℤ × ℤ /-- A region is a finite set of cells -/ abbrev Region := Finset Cell /- \#\# Prototile and Protoset -/ /- A prototile is the "shape" of a tile, represented as a list of cells with no duplicates. This design allows both: * Computation: iterate over...
\end{definition}

\begin{theorem}
  \label{translateRegion_card}
  \lean{translateRegion_card}
  \uses{Cell, Region, translateCell_injective, translateRegion}
  \leanok
  \proven
  A cell on the integer grid -/ abbrev Cell := ℤ × ℤ /-- A region is a finite set of cells -/ abbrev Region := Finset Cell /- \#\# Prototile and Protoset -/ /- A prototile is the "shape" of a tile, represented as a list of cells with no duplicates. This design allows both: * Computation: iterate over...
\end{theorem}

\begin{definition}
  \label{swapCell}
  \lean{swapCell}
  \uses{Cell}
  \leanok
  \proven
  A cell on the integer grid -/ abbrev Cell := ℤ × ℤ /-- A region is a finite set of cells -/ abbrev Region := Finset Cell /- \#\# Prototile and Protoset -/ /- A prototile is the "shape" of a tile, represented as a list of cells with no duplicates. This design allows both: * Computation: iterate over...
\end{definition}

\begin{theorem}
  \label{swapCell_involutive}
  \lean{swapCell_involutive}
  \uses{swapCell}
  \leanok
  \proven
  A cell on the integer grid -/ abbrev Cell := ℤ × ℤ /-- A region is a finite set of cells -/ abbrev Region := Finset Cell /- \#\# Prototile and Protoset -/ /- A prototile is the "shape" of a tile, represented as a list of cells with no duplicates. This design allows both: * Computation: iterate over...
\end{theorem}

\begin{theorem}
  \label{swapCell_injective}
  \lean{swapCell_injective}
  \uses{swapCell, swapCell_involutive}
  \leanok
  \proven
  A cell on the integer grid -/ abbrev Cell := ℤ × ℤ /-- A region is a finite set of cells -/ abbrev Region := Finset Cell /- \#\# Prototile and Protoset -/ /- A prototile is the "shape" of a tile, represented as a list of cells with no duplicates. This design allows both: * Computation: iterate over...
\end{theorem}

\begin{definition}
  \label{swapRegion}
  \lean{swapRegion}
  \uses{Region, swapCell}
  \leanok
  \proven
  A cell on the integer grid -/ abbrev Cell := ℤ × ℤ /-- A region is a finite set of cells -/ abbrev Region := Finset Cell /- \#\# Prototile and Protoset -/ /- A prototile is the "shape" of a tile, represented as a list of cells with no duplicates. This design allows both: * Computation: iterate over...
\end{definition}

\begin{definition}
  \label{PlacedTile_cells}
  \lean{PlacedTile.cells}
  \uses{Cell, PlacedTile, Protoset, rotateProto,
        translateCell}
  \leanok
  \proven
  A cell on the integer grid -/ abbrev Cell := ℤ × ℤ /-- A region is a finite set of cells -/ abbrev Region := Finset Cell /- \#\# Prototile and Protoset -/ /- A prototile is the "shape" of a tile, represented as a list of cells with no duplicates. This design allows both: * Computation: iterate over...
\end{definition}

\begin{definition}
  \label{PlacedTile_translate}
  \lean{PlacedTile.translate}
  \uses{Cell, PlacedTile, translate}
  \leanok
  \proven
  A cell on the integer grid -/ abbrev Cell := ℤ × ℤ /-- A region is a finite set of cells -/ abbrev Region := Finset Cell /- \#\# Prototile and Protoset -/ /- A prototile is the "shape" of a tile, represented as a list of cells with no duplicates. This design allows both: * Computation: iterate over...
\end{definition}

\begin{theorem}
  \label{PlacedTile_translate_cells}
  \lean{PlacedTile.translate_cells}
  \uses{Cell, PlacedTile, PlacedTile_cells, PlacedTile_translate,
        Protoset, translate, translateCell, translateRegion}
  \leanok
  \proven
  A cell on the integer grid -/ abbrev Cell := ℤ × ℤ /-- A region is a finite set of cells -/ abbrev Region := Finset Cell /- \#\# Prototile and Protoset -/ /- A prototile is the "shape" of a tile, represented as a list of cells with no duplicates. This design allows both: * Computation: iterate over...
\end{theorem}

\begin{theorem}
  \label{PlacedTile_cells_card}
  \lean{PlacedTile.cells_card}
  \uses{PlacedTile, PlacedTile_cells, Protoset, TileSet,
        rotateProto_card, translateCell_injective}
  \leanok
  \proven
  A cell on the integer grid -/ abbrev Cell := ℤ × ℤ /-- A region is a finite set of cells -/ abbrev Region := Finset Cell /- \#\# Prototile and Protoset -/ /- A prototile is the "shape" of a tile, represented as a list of cells with no duplicates. This design allows both: * Computation: iterate over...
\end{theorem}

\section{TileSet: Indexed Family of Placed Tiles}

\begin{definition}
  \label{TileSet}
  \lean{TileSet}
  \uses{PlacedTile, Protoset}
  \leanok
  \proven
  A cell on the integer grid -/ abbrev Cell := ℤ × ℤ /-- A region is a finite set of cells -/ abbrev Region := Finset Cell /- \#\# Prototile and Protoset -/ /- A prototile is the "shape" of a tile, represented as a list of cells with no duplicates. This design allows both: * Computation: iterate over...
\end{definition}

\begin{definition}
  \label{cellsAt}
  \lean{cellsAt}
  \uses{Cell, TileSet}
  \leanok
  \proven
  A cell on the integer grid -/ abbrev Cell := ℤ × ℤ /-- A region is a finite set of cells -/ abbrev Region := Finset Cell /- \#\# Prototile and Protoset -/ /- A prototile is the "shape" of a tile, represented as a list of cells with no duplicates. This design allows both: * Computation: iterate over...
\end{definition}

\begin{definition}
  \label{coveredCells}
  \lean{coveredCells}
  \uses{Cell, TileSet, cellsAt}
  \leanok
  \proven
  A cell on the integer grid -/ abbrev Cell := ℤ × ℤ /-- A region is a finite set of cells -/ abbrev Region := Finset Cell /- \#\# Prototile and Protoset -/ /- A prototile is the "shape" of a tile, represented as a list of cells with no duplicates. This design allows both: * Computation: iterate over...
\end{definition}

\begin{definition}
  \label{PairwiseDisjoint}
  \lean{PairwiseDisjoint}
  \uses{TileSet, cellsAt}
  \leanok
  \proven
  A cell on the integer grid -/ abbrev Cell := ℤ × ℤ /-- A region is a finite set of cells -/ abbrev Region := Finset Cell /- \#\# Prototile and Protoset -/ /- A prototile is the "shape" of a tile, represented as a list of cells with no duplicates. This design allows both: * Computation: iterate over...
\end{definition}

\begin{theorem}
  \label{card_coveredCells}
  \lean{card_coveredCells}
  \uses{PairwiseDisjoint, TileSet, cellsAt, coveredCells}
  \leanok
  \proven
  A cell on the integer grid -/ abbrev Cell := ℤ × ℤ /-- A region is a finite set of cells -/ abbrev Region := Finset Cell /- \#\# Prototile and Protoset -/ /- A prototile is the "shape" of a tile, represented as a list of cells with no duplicates. This design allows both: * Computation: iterate over...
\end{theorem}

\begin{definition}
  \label{Valid}
  \lean{Valid}
  \uses{PairwiseDisjoint, Region, TileSet, coveredCells}
  \leanok
  \proven
  A cell on the integer grid -/ abbrev Cell := ℤ × ℤ /-- A region is a finite set of cells -/ abbrev Region := Finset Cell /- \#\# Prototile and Protoset -/ /- A prototile is the "shape" of a tile, represented as a list of cells with no duplicates. This design allows both: * Computation: iterate over...
\end{definition}

\begin{definition}
  \label{instDecidablePairwiseDisjoint}
  \lean{instDecidablePairwiseDisjoint}
  \uses{PairwiseDisjoint, TileSet, cellsAt}
  \leanok
  \proven
  A cell on the integer grid -/ abbrev Cell := ℤ × ℤ /-- A region is a finite set of cells -/ abbrev Region := Finset Cell /- \#\# Prototile and Protoset -/ /- A prototile is the "shape" of a tile, represented as a list of cells with no duplicates. This design allows both: * Computation: iterate over...
\end{definition}

\begin{definition}
  \label{instDecidableValid}
  \lean{instDecidableValid}
  \uses{PairwiseDisjoint, Region, TileSet, Valid,
        coveredCells}
  \leanok
  \proven
  A cell on the integer grid -/ abbrev Cell := ℤ × ℤ /-- A region is a finite set of cells -/ abbrev Region := Finset Cell /- \#\# Prototile and Protoset -/ /- A prototile is the "shape" of a tile, represented as a list of cells with no duplicates. This design allows both: * Computation: iterate over...
\end{definition}

\begin{theorem}
  \label{Valid_disjoint_tiles}
  \lean{Valid.disjoint_tiles}
  \uses{Region, TileSet, Valid, cellsAt}
  \leanok
  \proven
  A cell on the integer grid -/ abbrev Cell := ℤ × ℤ /-- A region is a finite set of cells -/ abbrev Region := Finset Cell /- \#\# Prototile and Protoset -/ /- A prototile is the "shape" of a tile, represented as a list of cells with no duplicates. This design allows both: * Computation: iterate over...
\end{theorem}

\begin{theorem}
  \label{Valid_tile_contained}
  \lean{Valid.tile_contained}
  \uses{Region, TileSet, Valid, cellsAt,
        coveredCells}
  \leanok
  \proven
  A cell on the integer grid -/ abbrev Cell := ℤ × ℤ /-- A region is a finite set of cells -/ abbrev Region := Finset Cell /- \#\# Prototile and Protoset -/ /- A prototile is the "shape" of a tile, represented as a list of cells with no duplicates. This design allows both: * Computation: iterate over...
\end{theorem}

\begin{definition}
  \label{translate}
  \lean{translate}
  \uses{Cell, TileSet}
  \leanok
  \proven
  A cell on the integer grid -/ abbrev Cell := ℤ × ℤ /-- A region is a finite set of cells -/ abbrev Region := Finset Cell /- \#\# Prototile and Protoset -/ /- A prototile is the "shape" of a tile, represented as a list of cells with no duplicates. This design allows both: * Computation: iterate over...
\end{definition}

\begin{theorem}
  \label{translate_preserves_disjoint}
  \lean{translate_preserves_disjoint}
  \uses{Cell, PairwiseDisjoint, PlacedTile, PlacedTile_translate_cells,
        TileSet, cellsAt, translate, translateCell,
        translateCell_injective, translateRegion}
  \leanok
  \proven
  A cell on the integer grid -/ abbrev Cell := ℤ × ℤ /-- A region is a finite set of cells -/ abbrev Region := Finset Cell /- \#\# Prototile and Protoset -/ /- A prototile is the "shape" of a tile, represented as a list of cells with no duplicates. This design allows both: * Computation: iterate over...
\end{theorem}

\begin{theorem}
  \label{translate_coveredCells}
  \lean{translate_coveredCells}
  \uses{Cell, PlacedTile, PlacedTile_translate_cells, TileSet,
        coveredCells, translate, translateRegion}
  \leanok
  \proven
  A cell on the integer grid -/ abbrev Cell := ℤ × ℤ /-- A region is a finite set of cells -/ abbrev Region := Finset Cell /- \#\# Prototile and Protoset -/ /- A prototile is the "shape" of a tile, represented as a list of cells with no duplicates. This design allows both: * Computation: iterate over...
\end{theorem}

\begin{theorem}
  \label{translate_preserves_valid}
  \lean{translate_preserves_valid}
  \uses{Cell, Region, TileSet, Valid,
        translate, translateRegion, translate_coveredCells, translate_preserves_disjoint}
  \leanok
  \proven
  A cell on the integer grid -/ abbrev Cell := ℤ × ℤ /-- A region is a finite set of cells -/ abbrev Region := Finset Cell /- \#\# Prototile and Protoset -/ /- A prototile is the "shape" of a tile, represented as a list of cells with no duplicates. This design allows both: * Computation: iterate over...
\end{theorem}

\section{General Tileability}

\begin{definition}
  \label{Tileable}
  \lean{Tileable}
  \uses{Protoset, Region, TileSet, Valid}
  \leanok
  \proven
  A cell on the integer grid -/ abbrev Cell := ℤ × ℤ /-- A region is a finite set of cells -/ abbrev Region := Finset Cell /- \#\# Prototile and Protoset -/ /- A prototile is the "shape" of a tile, represented as a list of cells with no duplicates. This design allows both: * Computation: iterate over...
\end{definition}

\begin{theorem}
  \label{empty_tileable}
  \lean{empty_tileable}
  \uses{Protoset, TileSet, Tileable, coveredCells}
  \leanok
  \proven
  A cell on the integer grid -/ abbrev Cell := ℤ × ℤ /-- A region is a finite set of cells -/ abbrev Region := Finset Cell /- \#\# Prototile and Protoset -/ /- A prototile is the "shape" of a tile, represented as a list of cells with no duplicates. This design allows both: * Computation: iterate over...
\end{theorem}

\begin{theorem}
  \label{Tileable_translate}
  \lean{Tileable_translate}
  \uses{Cell, Protoset, Region, Tileable,
        translate, translateRegion, translate_preserves_valid}
  \leanok
  \proven
  A cell on the integer grid -/ abbrev Cell := ℤ × ℤ /-- A region is a finite set of cells -/ abbrev Region := Finset Cell /- \#\# Prototile and Protoset -/ /- A prototile is the "shape" of a tile, represented as a list of cells with no duplicates. This design allows both: * Computation: iterate over...
\end{theorem}

\begin{theorem}
  \label{Tileable_union}
  \lean{Tileable_union}
  \uses{Protoset, Region, TileSet, Tileable,
        cellsAt, coveredCells}
  \leanok
  \proven
  A cell on the integer grid -/ abbrev Cell := ℤ × ℤ /-- A region is a finite set of cells -/ abbrev Region := Finset Cell /- \#\# Prototile and Protoset -/ /- A prototile is the "shape" of a tile, represented as a list of cells with no duplicates. This design allows both: * Computation: iterate over...
\end{theorem}

\begin{theorem}
  \label{Tileable_biUnion}
  \lean{Tileable_biUnion}
  \uses{Protoset, Region, Tileable, Tileable_union,
        empty_tileable}
  \leanok
  \proven
  A cell on the integer grid -/ abbrev Cell := ℤ × ℤ /-- A region is a finite set of cells -/ abbrev Region := Finset Cell /- \#\# Prototile and Protoset -/ /- A prototile is the "shape" of a tile, represented as a list of cells with no duplicates. This design allows both: * Computation: iterate over...
\end{theorem}

\begin{theorem}
  \label{translateRegion_neg}
  \lean{translateRegion_neg}
  \uses{Cell, Region, translateCell, translateRegion}
  \leanok
  \proven
  A cell on the integer grid -/ abbrev Cell := ℤ × ℤ /-- A region is a finite set of cells -/ abbrev Region := Finset Cell /- \#\# Prototile and Protoset -/ /- A prototile is the "shape" of a tile, represented as a list of cells with no duplicates. This design allows both: * Computation: iterate over...
\end{theorem}

\begin{theorem}
  \label{translateRegion_add}
  \lean{translateRegion_add}
  \uses{Cell, Region, translateCell, translateRegion}
  \leanok
  \proven
  A cell on the integer grid -/ abbrev Cell := ℤ × ℤ /-- A region is a finite set of cells -/ abbrev Region := Finset Cell /- \#\# Prototile and Protoset -/ /- A prototile is the "shape" of a tile, represented as a list of cells with no duplicates. This design allows both: * Computation: iterate over...
\end{theorem}

\begin{theorem}
  \label{erase_union_finset}
  \lean{erase_union_finset}
  \leanok
  \proven
  A cell on the integer grid -/ abbrev Cell := ℤ × ℤ /-- A region is a finite set of cells -/ abbrev Region := Finset Cell /- \#\# Prototile and Protoset -/ /- A prototile is the "shape" of a tile, represented as a list of cells with no duplicates. This design allows both: * Computation: iterate over...
\end{theorem}

\begin{theorem}
  \label{translateRegion_erase_point}
  \lean{translateRegion_erase_point}
  \uses{Cell, Region, translateCell, translateCell_injective,
        translateRegion}
  \leanok
  \proven
  A cell on the integer grid -/ abbrev Cell := ℤ × ℤ /-- A region is a finite set of cells -/ abbrev Region := Finset Cell /- \#\# Prototile and Protoset -/ /- A prototile is the "shape" of a tile, represented as a list of cells with no duplicates. This design allows both: * Computation: iterate over...
\end{theorem}

\begin{theorem}
  \label{Tileable_translate_iff}
  \lean{Tileable_translate_iff}
  \uses{Cell, Protoset, Region, Tileable,
        Tileable_translate, translateRegion, translateRegion_neg}
  \leanok
  \proven
  A cell on the integer grid -/ abbrev Cell := ℤ × ℤ /-- A region is a finite set of cells -/ abbrev Region := Finset Cell /- \#\# Prototile and Protoset -/ /- A prototile is the "shape" of a tile, represented as a list of cells with no duplicates. This design allows both: * Computation: iterate over...
\end{theorem}

\begin{theorem}
  \label{Tileable_of_translate_eq}
  \lean{Tileable_of_translate_eq}
  \uses{Cell, Protoset, Region, TileSet,
        Tileable, Tileable_translate_iff, translateRegion}
  \leanok
  \proven
  A cell on the integer grid -/ abbrev Cell := ℤ × ℤ /-- A region is a finite set of cells -/ abbrev Region := Finset Cell /- \#\# Prototile and Protoset -/ /- A prototile is the "shape" of a tile, represented as a list of cells with no duplicates. This design allows both: * Computation: iterate over...
\end{theorem}

\section{Helper for Creating TileSets}

\begin{definition}
  \label{mkTileSet}
  \lean{mkTileSet}
  \uses{PlacedTile, Protoset, TileSet}
  \leanok
  \proven
  A cell on the integer grid -/ abbrev Cell := ℤ × ℤ /-- A region is a finite set of cells -/ abbrev Region := Finset Cell /- \#\# Prototile and Protoset -/ /- A prototile is the "shape" of a tile, represented as a list of cells with no duplicates. This design allows both: * Computation: iterate over...
\end{definition}

\begin{definition}
  \label{mkPlacedTile}
  \lean{mkPlacedTile}
  \uses{PlacedTile}
  \leanok
  \proven
  A cell on the integer grid -/ abbrev Cell := ℤ × ℤ /-- A region is a finite set of cells -/ abbrev Region := Finset Cell /- \#\# Prototile and Protoset -/ /- A prototile is the "shape" of a tile, represented as a list of cells with no duplicates. This design allows both: * Computation: iterate over...
\end{definition}

\section{Generic Placement Enumeration}

\begin{definition}
  \label{PlacedTile_containedIn}
  \lean{PlacedTile.containedIn}
  \uses{PlacedTile, Protoset, Region}
  \leanok
  \proven
  A cell on the integer grid -/ abbrev Cell := ℤ × ℤ /-- A region is a finite set of cells -/ abbrev Region := Finset Cell /- \#\# Prototile and Protoset -/ /- A prototile is the "shape" of a tile, represented as a list of cells with no duplicates. This design allows both: * Computation: iterate over...
\end{definition}

\begin{definition}
  \label{PlacedTile_coversCell}
  \lean{PlacedTile.coversCell}
  \uses{Cell, PlacedTile, Protoset, rotateCell}
  \leanok
  \proven
  A cell on the integer grid -/ abbrev Cell := ℤ × ℤ /-- A region is a finite set of cells -/ abbrev Region := Finset Cell /- \#\# Prototile and Protoset -/ /- A prototile is the "shape" of a tile, represented as a list of cells with no duplicates. This design allows both: * Computation: iterate over...
\end{definition}

\begin{definition}
  \label{allPlacementsCoveringProto}
  \lean{allPlacementsCoveringProto}
  \uses{Cell, PlacedTile, Prototile, rotateCell}
  \leanok
  \proven
  A cell on the integer grid -/ abbrev Cell := ℤ × ℤ /-- A region is a finite set of cells -/ abbrev Region := Finset Cell /- \#\# Prototile and Protoset -/ /- A prototile is the "shape" of a tile, represented as a list of cells with no duplicates. This design allows both: * Computation: iterate over...
\end{definition}

\begin{definition}
  \label{allPlacementsCovering}
  \lean{allPlacementsCovering}
  \uses{Cell, PlacedTile, Protoset, allPlacementsCoveringProto}
  \leanok
  \proven
  A cell on the integer grid -/ abbrev Cell := ℤ × ℤ /-- A region is a finite set of cells -/ abbrev Region := Finset Cell /- \#\# Prototile and Protoset -/ /- A prototile is the "shape" of a tile, represented as a list of cells with no duplicates. This design allows both: * Computation: iterate over...
\end{definition}

\begin{theorem}
  \label{allPlacementsCoveringProto_covers}
  \lean{allPlacementsCoveringProto_covers}
  \uses{Cell, PlacedTile, PlacedTile_cells, PlacedTile_coversCell,
        Protoset, Prototile, Prototile_mem_coe, allPlacementsCoveringProto,
        rotateCell, rotateProto, translateCell}
  \leanok
  \proven
  A cell on the integer grid -/ abbrev Cell := ℤ × ℤ /-- A region is a finite set of cells -/ abbrev Region := Finset Cell /- \#\# Prototile and Protoset -/ /- A prototile is the "shape" of a tile, represented as a list of cells with no duplicates. This design allows both: * Computation: iterate over...
\end{theorem}

\begin{theorem}
  \label{allPlacementsCoveringProto_complete}
  \lean{allPlacementsCoveringProto_complete}
  \uses{Cell, PlacedTile, PlacedTile_cells, PlacedTile_coversCell,
        Protoset, allPlacementsCoveringProto, rotateProto, translateCell}
  \leanok
  \proven
  A cell on the integer grid -/ abbrev Cell := ℤ × ℤ /-- A region is a finite set of cells -/ abbrev Region := Finset Cell /- \#\# Prototile and Protoset -/ /- A prototile is the "shape" of a tile, represented as a list of cells with no duplicates. This design allows both: * Computation: iterate over...
\end{theorem}

\begin{definition}
  \label{placementsCoveringIn}
  \lean{placementsCoveringIn}
  \uses{Cell, PlacedTile, Protoset, Region,
        allPlacementsCovering}
  \leanok
  \proven
  A cell on the integer grid -/ abbrev Cell := ℤ × ℤ /-- A region is a finite set of cells -/ abbrev Region := Finset Cell /- \#\# Prototile and Protoset -/ /- A prototile is the "shape" of a tile, represented as a list of cells with no duplicates. This design allows both: * Computation: iterate over...
\end{definition}

\section{Protoset Refinement}

\begin{definition}
  \label{Refines}
  \lean{Refines}
  \uses{Protoset, Tileable, rotateProto}
  \leanok
  \proven
  A cell on the integer grid -/ abbrev Cell := ℤ × ℤ /-- A region is a finite set of cells -/ abbrev Region := Finset Cell /- \#\# Prototile and Protoset -/ /- A prototile is the "shape" of a tile, represented as a list of cells with no duplicates. This design allows both: * Computation: iterate over...
\end{definition}

\begin{theorem}
  \label{Tileable_refine}
  \lean{Tileable_refine}
  \uses{PlacedTile, PlacedTile_cells, Protoset, Refines,
        Region, TileSet, Tileable, Tileable_biUnion,
        Tileable_translate, cellsAt, coveredCells, rectangle,
        rotateProto, translateRegion}
  \leanok
  \proven
  A cell on the integer grid -/ abbrev Cell := ℤ × ℤ /-- A region is a finite set of cells -/ abbrev Region := Finset Cell /- \#\# Prototile and Protoset -/ /- A prototile is the "shape" of a tile, represented as a list of cells with no duplicates. This design allows both: * Computation: iterate over...
\end{theorem}

\section{Rectangle Definition and Tilings}

\begin{definition}
  \label{rectangle}
  \lean{rectangle}
  \uses{Region}
  \leanok
  \proven
  A cell on the integer grid -/ abbrev Cell := ℤ × ℤ /-- A region is a finite set of cells -/ abbrev Region := Finset Cell /- \#\# Prototile and Protoset -/ /- A prototile is the "shape" of a tile, represented as a list of cells with no duplicates. This design allows both: * Computation: iterate over...
\end{definition}

\begin{theorem}
  \label{mem_rectangle}
  \lean{mem_rectangle}
  \uses{Cell, rectangle}
  \leanok
  \proven
  A cell on the integer grid -/ abbrev Cell := ℤ × ℤ /-- A region is a finite set of cells -/ abbrev Region := Finset Cell /- \#\# Prototile and Protoset -/ /- A prototile is the "shape" of a tile, represented as a list of cells with no duplicates. This design allows both: * Computation: iterate over...
\end{theorem}

\begin{theorem}
  \label{swap_rectangle}
  \lean{swap_rectangle}
  \uses{rectangle, swapCell, swapRegion}
  \leanok
  \proven
  A cell on the integer grid -/ abbrev Cell := ℤ × ℤ /-- A region is a finite set of cells -/ abbrev Region := Finset Cell /- \#\# Prototile and Protoset -/ /- A prototile is the "shape" of a tile, represented as a list of cells with no duplicates. This design allows both: * Computation: iterate over...
\end{theorem}

\section{Rectangle Splitting and Combining}

\begin{theorem}
  \label{rectangle_split_horizontal}
  \lean{rectangle_split_horizontal}
  \uses{mem_rectangle, rectangle, translateCell, translateRegion}
  \leanok
  \proven
  A cell on the integer grid -/ abbrev Cell := ℤ × ℤ /-- A region is a finite set of cells -/ abbrev Region := Finset Cell /- \#\# Prototile and Protoset -/ /- A prototile is the "shape" of a tile, represented as a list of cells with no duplicates. This design allows both: * Computation: iterate over...
\end{theorem}

\begin{theorem}
  \label{rectangle_split_horizontal_disjoint}
  \lean{rectangle_split_horizontal_disjoint}
  \uses{mem_rectangle, rectangle, translateCell, translateRegion}
  \leanok
  \proven
  A cell on the integer grid -/ abbrev Cell := ℤ × ℤ /-- A region is a finite set of cells -/ abbrev Region := Finset Cell /- \#\# Prototile and Protoset -/ /- A prototile is the "shape" of a tile, represented as a list of cells with no duplicates. This design allows both: * Computation: iterate over...
\end{theorem}

\begin{theorem}
  \label{rectangle_split_vertical}
  \lean{rectangle_split_vertical}
  \uses{mem_rectangle, rectangle, translateCell, translateRegion}
  \leanok
  \proven
  A cell on the integer grid -/ abbrev Cell := ℤ × ℤ /-- A region is a finite set of cells -/ abbrev Region := Finset Cell /- \#\# Prototile and Protoset -/ /- A prototile is the "shape" of a tile, represented as a list of cells with no duplicates. This design allows both: * Computation: iterate over...
\end{theorem}

\begin{theorem}
  \label{rectangle_split_vertical_disjoint}
  \lean{rectangle_split_vertical_disjoint}
  \uses{mem_rectangle, rectangle, translateCell, translateRegion}
  \leanok
  \proven
  A cell on the integer grid -/ abbrev Cell := ℤ × ℤ /-- A region is a finite set of cells -/ abbrev Region := Finset Cell /- \#\# Prototile and Protoset -/ /- A prototile is the "shape" of a tile, represented as a list of cells with no duplicates. This design allows both: * Computation: iterate over...
\end{theorem}

\begin{theorem}
  \label{Tileable_horizontal_union}
  \lean{Tileable_horizontal_union}
  \uses{Protoset, Tileable, Tileable_translate, Tileable_union,
        rectangle, rectangle_split_horizontal, rectangle_split_horizontal_disjoint}
  \leanok
  \proven
  A cell on the integer grid -/ abbrev Cell := ℤ × ℤ /-- A region is a finite set of cells -/ abbrev Region := Finset Cell /- \#\# Prototile and Protoset -/ /- A prototile is the "shape" of a tile, represented as a list of cells with no duplicates. This design allows both: * Computation: iterate over...
\end{theorem}

\begin{theorem}
  \label{Tileable_vertical_union}
  \lean{Tileable_vertical_union}
  \uses{Protoset, Tileable, Tileable_translate, Tileable_union,
        rectangle, rectangle_split_vertical, rectangle_split_vertical_disjoint}
  \leanok
  \proven
  A cell on the integer grid -/ abbrev Cell := ℤ × ℤ /-- A region is a finite set of cells -/ abbrev Region := Finset Cell /- \#\# Prototile and Protoset -/ /- A prototile is the "shape" of a tile, represented as a list of cells with no duplicates. This design allows both: * Computation: iterate over...
\end{theorem}

